

\chapter{抹茶题集}
\begin{example}[新曲线]
已知曲线 $\Gamma: 12x^3 + 12xy^2 + 3x^2 + 4y^2 = 0$。过原点 $O$ 作两条互相垂直的直线，分别交曲线 $\Gamma$ 于异于原点的两点 $A, B$。
求证：无论直线如何转动，$\triangle OAB$ 的外接圆恒过一个异于原点的定点，并求出该定点坐标。
\end{example}
\begin{solution}
设其中一条直线斜率为 $m\neq 0$，另一条直线斜率为 $-\frac{1}{m}$；若有一条直线是竖直的，交换两线即可。

设直线 $OA$ 的方程为 $y=mx$，则其与曲线 $\Gamma$ 的非原点交点 $A$ 满足
\[
12x^3+12m^2x^3+3x^2+4m^2x^2=0.
\]
由此得
\[
x_A=-\frac{3+4m^2}{12(1+m^2)},\qquad y_A=mx_A.
\]
同理，直线 $OB$ 的方程为 $y=-\frac{x}{m}$，其与曲线 $\Gamma$ 的非原点交点 $B$ 满足
\[
x_B=-\frac{3m^2+4}{12(1+m^2)},\qquad y_B=-\frac{x_B}{m}.
\]

设过 $O,A,B$ 的圆方程为
\[
x^2+y^2+Dx+Ey=0.
\]
代入点 $A,B$，得
\[
D+ mE=\frac{3+4m^2}{12},\qquad D-\frac{E}{m}=\frac{3m^2+4}{12m^2}.
\]
第一式保持不变，第二式乘以 $m^2$ 后相加，消去 $E$，得到
\[
(m^2+1)D=\frac{7(m^2+1)}{12},
\]
故
\[
D=\frac{7}{12}.
\]
于是圆方程可写成
\[
x^2+y^2+\frac{7}{12}x+Ey=0.
\]
令 $y=0$，得到 $x\left(x+\frac{7}{12}\right)=0$，故该圆恒过定点
\[
\left(-\frac{7}{12},0\right).
\]
\end{solution}
\begin{example}
（多选）已知非零复数 $z_1, z_2$ 满足 $z\bar{z} = z + \bar{z}$。设复数 $z_3$ 满足$\frac{2}{z_3} = \frac{1}{z_1} + \frac{1}{z_2}$。下列关于 $z_1, z_2, z_3$ 的结论中，正确的有（~~~~~）\\
A. 复数 $z_3$ 必定满足方程 $z\bar{z} = z + \bar{z}$\\B. $\frac{|z_1 - z_3|}{|z_1|} = \frac{|z_2 - z_3|}{|z_2|}$\\C. $\frac{1}{|z_1|} + \frac{1}{|z_2|} \ge \frac{2}{|z_3|}$\\D. $\frac{1}{|z_1|^2} + \frac{1}{|z_2|^2} \ge \frac{2}{|z_3|^2}$
\end{example}
\begin{solution}
对 A，由
\[
\frac{2}{z_3}=\frac{1}{z_1}+\frac{1}{z_2}
\]
可得
\[
z_3=\frac{2z_1z_2}{z_1+z_2},\qquad \bar z_3=\frac{2\bar z_1\bar z_2}{\bar z_1+\bar z_2}.
\]
于是
\[
z_3\bar z_3=\frac{4(z_1\bar z_1)(z_2\bar z_2)}{(z_1+z_2)(\bar z_1+\bar z_2)}.
\]
由已知条件 $z_1\bar z_1=z_1+\bar z_1,\ z_2\bar z_2=z_2+\bar z_2$，得
\[
z_3\bar z_3=\frac{4(z_1+\bar z_1)(z_2+\bar z_2)}{(z_1+z_2)(\bar z_1+\bar z_2)}.
\]
另一方面，
\[
z_3+\bar z_3=\frac{2z_1z_2}{z_1+z_2}+\frac{2\bar z_1\bar z_2}{\bar z_1+\bar z_2}.
\]
通分后得
\[
z_3+\bar z_3
=\frac{2z_1z_2(\bar z_1+\bar z_2)+2\bar z_1\bar z_2(z_1+z_2)}{(z_1+z_2)(\bar z_1+\bar z_2)}.
\]
分子可整理为
\[
2(z_1\bar z_1)z_2+2(z_2\bar z_2)z_1+2(z_1\bar z_1)\bar z_2+2(z_2\bar z_2)\bar z_1.
\]
再代入 $z\bar z=z+\bar z$，可得
\[
z_3+\bar z_3=\frac{4(z_1+\bar z_1)(z_2+\bar z_2)}{(z_1+z_2)(\bar z_1+\bar z_2)}.
\]
因此 $z_3\bar z_3=z_3+\bar z_3$，A 正确。

对 B，由
\[
z_1-z_3=z_1-\frac{2z_1z_2}{z_1+z_2}
=\frac{z_1(z_1-z_2)}{z_1+z_2},
\]
所以
\[
\frac{|z_1-z_3|}{|z_1|}
=\frac{|z_1-z_2|}{|z_1+z_2|}.
\]
同理，
\[
z_2-z_3=z_2-\frac{2z_1z_2}{z_1+z_2}
=\frac{z_2(z_2-z_1)}{z_1+z_2},
\]
从而
\[
\frac{|z_2-z_3|}{|z_2|}
=\frac{|z_2-z_1|}{|z_1+z_2|}
=\frac{|z_1-z_2|}{|z_1+z_2|}.
\]
故
\[
\frac{|z_1-z_3|}{|z_1|}=\frac{|z_2-z_3|}{|z_2|},
\]
所以 B 正确。

对 C、D，令
\[
w_i=\frac1{z_i}\quad(i=1,2,3).
\]
由 $z\bar z=z+\bar z$ 可得
\[
w+\bar w=1,
\]
故 $w_1,w_2$ 在同一直线 $\operatorname{Re}w=\frac12$ 上。又由题设
\[
w_3=\frac{w_1+w_2}{2}.
\]
于是
\[
|w_1|+|w_2|\ge |w_1+w_2|=2|w_3|,
\]
还原即为 C 正确。并且
\[
|w_1|^2+|w_2|^2
=2|w_3|^2+\frac12|w_1-w_2|^2\ge2|w_3|^2,
\]
还原即为 D 正确。
\end{solution}
\begin{example}[抹茶T1]
    椭圆 $\frac{x^2}{4}+y^2=m\ (m>1)$ 上三点 $A,B,P(0,1)$满足$\overrightarrow{AP}=2\overrightarrow{PB}$，求 $|x_B|$ 的最大值
\end{example}
\begin{solution}
由 $\overrightarrow{AP}=2\overrightarrow{PB}$ 可以把点 $A$ 用点 $B$ 表示，题目所求随即转化为点 $B$ 轨迹上 $|x_B|$ 的上界。设点 $A(x_A, y_A)$，点 $B(x_B, y_B)$。已知定点 $P(0,1)$。有
$$ \overrightarrow{AP} = (0 - x_A, 1 - y_A) = (-x_A, 1 - y_A), \overrightarrow{PB} = (x_B - 0, y_B - 1) = (x_B, y_B - 1) $$
由于 $\overrightarrow{AP} = 2\overrightarrow{PB}$，代入条件：$$ \begin{cases} -x_A = 2x_B \\ 1 - y_A = 2(y_B - 1) \end{cases} \implies \begin{cases} \mathbf{x_A = -2x_B} \\ \mathbf{y_A = 3 - 2y_B} \end{cases} $$
已知 $A$ 和 $B$ 都在椭圆 $\frac{x^2}{4} + y^2 = m$ 上。分别代入
$$ \frac{x_B^2}{4} + y_B^2 = m ,\frac{(-2x_B)^2}{4} + (3 - 2y_B)^2 = m $$
消去 $x_B$得到：$$ \frac{x_B^2}{4} + y_B^2 = x_B^2 + 4y_B^2 - 12y_B + 9 \Leftrightarrow\left(1 - \frac{1}{4}\right)x_B^2 + (4 - 1)y_B^2 - 12y_B + 9 = 0 $$
整理得到点 $B$ 的轨迹
\[
E_B:\ \frac{x^2}{4}+(y-2)^2=1.
\]
由
\[
m=\frac{x_B^2}{4}+y_B^2=4y_B-3
\]
可知，条件 $m>1$ 等价于
\[
4y_B-3>1 \iff y_B>1.
\]
因此，满足题意的点 $B$ 恰好是椭圆
\[
\frac{x_B^2}{4}+(y_B-2)^2=1
\]
上满足 $y_B>1$ 的点。

又因为
\[
\frac{x_B^2}{4}=1-(y_B-2)^2\le 1,
\]
所以
\[
x_B^2\le 4,\qquad |x_B|\le 2.
\]
当且仅当 $y_B=2$ 时取等号，这时 $m=4\cdot 2-3=5>1$，故
\[
|x_B|_{\max}=2.
\]
\end{solution}
\begin{example}[抹茶T2]
    椭圆 $\frac{x^{2}}{4}+y^{2}=1$ 的右焦点为 $F$，过点 $F$ 的直线交椭圆于 $A,B$ 两点，若 $|OA|^{2}=|AB|$，求点 $A$ 的坐标.
\end{example}
\begin{solution}
    设点 $A,B$ 对应的焦半径分别为 $r_A,r_B$。对于椭圆上过焦点的弦，有
    \[
    \frac{1}{r_A}+\frac{1}{r_B}=\frac{2a}{b^2}.
    \]
    这里 $a=2,\ b=1$，所以
    \[
    \frac{1}{r_A}+\frac{1}{r_B}=4.
    \]
    又
    \[
    r_A=2-\frac{\sqrt3}{2}x_A,
    \]
    从而
    \[
    \frac{1}{r_B}=4-\frac{1}{r_A}=\frac{4r_A-1}{r_A},
    \qquad
    r_B=\frac{r_A}{4r_A-1}.
    \]
    因为 $F$ 在弦 $AB$ 上，且位于 $A,B$ 之间，所以
    \[
    |AB|=r_A+r_B=r_A+\frac{r_A}{4r_A-1}=\frac{4r_A^2}{4r_A-1}.
    \]
    将 $r_A = 2 - \frac{\sqrt{3}}{2}x_A$ 代入，得
    \[
    |AB|=\frac{3x_A^2 - 8\sqrt{3}x_A + 16}{7 - 2\sqrt{3}x_A}.
    \]
    另一方面，
    \[
    |OA|^2=x_A^2+y_A^2=x_A^2+1-\frac{x_A^2}{4}=\frac{3x_A^2+4}{4}.
    \]
    由题设 $|OA|^2=|AB|$，得
    \[
    2\sqrt{3}x_A^3-3x_A^2-8\sqrt{3}x_A+12=0.
    \]
    分解为
    \[
    (x_A^2-4)(2\sqrt{3}x_A-3)=0.
    \]
    因而
    \[
    x_A=\pm2,\qquad x_A=\frac{\sqrt3}{2}.
    \]
    对应的点分别为 $(2,0)$、$(-2,0)$，以及
    \[
    \left(\frac{\sqrt3}{2},\frac{\sqrt{13}}{4}\right),\quad
    \left(\frac{\sqrt3}{2},-\frac{\sqrt{13}}{4}\right).
    \]
\end{solution}

\begin{example}[抹茶T3]
    抛物线 $y^2=4x$ 的焦点为 $F$，过焦点的直线交抛物线于 $A,B$ 两点，点 $M$ 在抛物线上，满足 $MA\perp MB$，$FM\perp AB$，求 $\triangle ABM$ 的面积。
\end{example}
\begin{solution}
    假设 $AB \perp x$ 轴，则直线方程为 $x = 1$。解得 $y^2 = 4 \implies A(1, 2), B(1, -2)$，此时 $\overrightarrow{MA} \cdot \overrightarrow{MB} = -3 \neq 0$，与 $MA \perp MB$ 矛盾。因此设直线 $AB$ 的方程为 $x = ty + 1$（其中 $t \neq 0$），联立得
    \[
    y^2-4ty-4=0.
    \]
    记交点对应的纵坐标为 $y_1,y_2$，则
    \[
    y_1+y_2=4t,\qquad y_1y_2=-4.
    \]
    设 $M(x_0,y_0)$，且 $x_0=\frac{y_0^2}{4}$。由 $MA\perp MB$ 得
    \[
    (y_1-y_0)(y_2-y_0)\left[\frac{(y_1+y_0)(y_2+y_0)}{16}+1\right]=0.
    \]
    因 $M\neq A,B$，故
    \[
    (y_1+y_0)(y_2+y_0)=-16.
    \]
    代入 Vieta 关系，得
    \[
    y_0^2+4ty_0+12=0.
    \]
    又由 $FM\perp AB$ 知
    \[
    t(y_0^2-4)+4y_0=0.
    \]
    两式联立，得
    \[
    (y_0^2-12)(y_0^2+4)=0.
    \]
    因而 $y_0^2=12$，从而 $x_0=3$，$M(3,\pm2\sqrt3)$。再由第一式得 $t^2=3$。于是
    \[
    |FM|=\sqrt{(3-1)^2+12}=4.
    \]
    另一方面，
    \[
    |AB|^2=(t^2+1)(y_1-y_2)^2=(t^2+1)\big[(y_1+y_2)^2-4y_1y_2\big]=16(t^2+1)^2,
    \]
    所以 $|AB|=4(t^2+1)=16$。
    因而
    \[
    S_{\triangle ABM}=\frac12|AB|\cdot|FM|=32.
    \]
\end{solution}

\begin{example}[抹茶T4]
    点 $A(0,-1)$，点 $P,Q$ 满足点 $A,P,Q$ 共线，且 $|AP|\cdot |AQ|=3$，若 $k_{OQ}=3k_{OP}$，求点 $P$ 的轨迹方程.
\end{example}
\begin{solution}
若 $k_{OP}=0$，则 $k_{OQ}=3k_{OP}=0$，于是 $P,Q$ 都在 $x$ 轴上；但 $A(0,-1),P,Q$ 三点共线时，直线 $AP$ 与 $x$ 轴只交于点 $P$，所以只能有 $Q=P$，这与 $|AP|\cdot |AQ|=3$ 矛盾。零斜率情形已排除，以下设两条斜率均存在且非零。设动点 $P$ 的坐标为 $(x, y)$，点 $Q$ 的坐标为 $(x_Q, y_Q)$。此时点 $P$ 与点 $Q$ 的横坐标均满足 $x \neq 0$ 且 $x_Q \neq 0$。由斜率公式得 $\frac{y_Q}{x_Q} = \frac{3y}{x}$，即 $x y_Q = 3x_Q y$。

由于点 $A(0, -1)$ 与 $P, Q$ 三点共线，且 $P$ 不在 $y$ 轴上，故向量 $\overrightarrow{AP} \neq \vec{0}$。设 $\overrightarrow{AQ} = \lambda \overrightarrow{AP}$（$\lambda \neq 0$），则有：
$$ (x_Q, y_Q + 1) = \lambda (x, y + 1) \implies x_Q = \lambda x, \quad y_Q = \lambda(y + 1) - 1 $$
将上述关系代入 $x y_Q = 3x_Q y$ 中，得 $x[\lambda(y + 1) - 1] = 3\lambda x y$。因 $x \neq 0$，等式两边同除以 $x$ 并整理得 $\lambda(y + 1) - 1 = 3\lambda y$，即 $\lambda(1 - 2y) = 1$。若 $y = \frac{1}{2}$，方程无解，故 $y \neq \frac{1}{2}$，从而解得 $\lambda = \frac{1}{1 - 2y}$。

根据长度条件 $|AP| \cdot |AQ| = 3$，结合 $\overrightarrow{AQ} = \lambda \overrightarrow{AP}$ 可得 $|\lambda| \cdot |AP|^2 = 3$。代入距离公式与 $\lambda$ 的表达式，得到关于 $P$ 点坐标的方程：
$$ \left| \frac{1}{1 - 2y} \right| \cdot [x^2 + (y + 1)^2] = 3 \implies x^2 + (y + 1)^2 = 3|1 - 2y| $$
针对 $1 - 2y$ 的正负号进行分类讨论：

当 $1 - 2y > 0$，即 $y < \frac{1}{2}$ 时，方程化为 $x^2 + y^2 + 2y + 1 = 3 - 6y$，整理得 $x^2 + y^2 + 8y - 2 = 0$，配方得标准方程 $x^2 + (y + 4)^2 = 18$。该圆上点的纵坐标最大值为 $y_{max} = -4 + \sqrt{18} = -4 + 3\sqrt{2}$。由于 $3\sqrt{2} = \sqrt{18} < 4.5$，故 $y_{max} < 0.5$，即该轨迹圆上的所有点均满足 $y < \frac{1}{2}$ 的前提条件。

当 $1 - 2y < 0$，即 $y > \frac{1}{2}$ 时，方程化为 $x^2 + y^2 + 2y + 1 = -(3 - 6y)$，整理得 $x^2 + y^2 - 4y + 4 = 0$，即 $x^2 + (y - 2)^2 = 0$。此方程的唯一实数解为点 $(0, 2)$，但该点横坐标 $x = 0$，不满足斜率存在的初始前提，故应予以舍弃。

因此，动点 $P$ 的轨迹是以 $(0, -4)$ 为圆心，$\sqrt{18}$ 为半径的圆，且需除去与 $y$ 轴的交点。其轨迹方程为：
$$ x^2 + (y + 4)^2 = 18 \quad (x \neq 0) $$
或写成一般式：$x^2 + y^2 + 8y - 2 = 0 \quad (x \neq 0)$。
\end{solution}
\begin{example}[抹茶T5]
    双曲线 $\frac{x^{2}}{a^{2}} - \frac{y^{2}}{b^{2}} = 1$ 的右顶点为 $A$, $\overrightarrow{OB} = 2\overrightarrow{OA}$, 若双曲线上存在点 $P$, 使得 $\angle APB = \frac{\pi}{2}$, 求双曲线的离心率的取值范围.
\end{example}
\begin{solution}
设双曲线的标准方程为 $\frac{x^2}{a^2} - \frac{y^2}{b^2} = 1$ ($a>0, b>0$)，其右顶点为 $A(a, 0)$。由 $\overrightarrow{OB} = 2\overrightarrow{OA}$ 可得 $B$ 点坐标为 $(2a, 0)$。设双曲线上存在点 $P(x, y)$ 满足 $\angle APB = \frac{\pi}{2}$，则有 $\overrightarrow{PA} \cdot \overrightarrow{PB} = 0$ 且点 $P$ 不与 $A, B$ 重合，即 $y \neq 0$。

根据向量坐标运算，$\overrightarrow{PA} = (a - x, -y)$，$\overrightarrow{PB} = (2a - x, -y)$，由数量积为零得：
$$(x - a)(x - 2a) + y^2 = 0 \implies y^2 = -(x - a)(x - 2a)$$
由于 $y^2 > 0$，可得 $x$ 的取值范围满足 $a < x < 2a$。此条件同时限定了点 $P$ 必须位于双曲线的右支。

将点 $P$ 在双曲线上的关系 $y^2 = b^2 \left( \frac{x^2}{a^2} - 1 \right)$ 代入上述方程。利用离心率关系 $\frac{b^2}{a^2} = e^2 - 1$，得 $y^2 = (e^2 - 1)(x^2 - a^2)$。联立圆与双曲线的方程：
$$(e^2 - 1)(x - a)(x + a) = -(x - a)(x - 2a)$$
由于 $x > a$，等式两边可约去 $x - a$，得：
$$(e^2 - 1)(x + a) = -(x - 2a)$$
整理得 $e^2 x + e^2 a - x - a = -x + 2a$，化简可解得点 $P$ 的横坐标为：
$$x = \frac{3 - e^2}{e^2} a$$

要使满足题意的点 $P$ 存在，该横坐标 $x$ 必须满足范围 $a < x < 2a$。代入得：
$$a < \frac{3 - e^2}{e^2} a < 2a$$
由于 $a > 0$，不等式各边同除以 $a$ 得 $1 < \frac{3}{e^2} - 1 < 2$，即 $2 < \frac{3}{e^2} < 3$。由此解得：
$$1 < e^2 < \frac{3}{2}$$
考虑到双曲线离心率 $e > 1$，对不等式开平方得 $1 < e < \frac{\sqrt{6}}{2}$。

经过检验，当 $e \in (1, \frac{\sqrt{6}}{2})$ 时，所求横坐标 $x$ 对应的纵坐标 $y = \pm \sqrt{-(x-a)(x-2a)}$ 必为非零实数，保证了点 $P$ 存在且不与 $A, B$ 重合，满足 $\angle APB = \frac{\pi}{2}$。若 $e = \frac{\sqrt{6}}{2}$，则 $x = a$，此时点 $P$ 与顶点 $A$ 重合，角不存在。因此，双曲线离心率 $e$ 的取值范围是 $(1, \frac{\sqrt{6}}{2})$。
\end{solution}

\chapter{反演变换在解析几何中的应用}

本节采用反演 $P\mapsto P'$，满足 $OP\cdot OP'=R^2$。在圆锥曲线问题中，过原点的线、圆关系常可转化为代数曲线上的交点条件，具体计算从点线圆对应方程开始。

\section{反演变换的基本原理}

设反演中心为 $O$，反演半径为 $R$。平面上任意一点 $P$（异于点 $O$）的反演像记为 $P^{\prime}$，它满足 $\overrightarrow{OP}$ 与 $\overrightarrow{OP}^{\prime}$ 同向共线，且 $OP\cdot OP^\prime=R^2$。

反演变换具有以下基本性质：
\begin{itemize}
    \item 过反演中心的直线变为过反演中心的直线。
    \item 不过反演中心的直线变为过反演中心的圆。
    \item 过反演中心的圆变为不过反演中心的直线。
    \item 不过反演中心的圆变为不过反演中心的圆。
    \item 反演是可逆的双射变换。若曲线 $\Gamma$ 的反演像为 $\Gamma^{\prime}$，则对 $\Gamma^{\prime}$ 再作同一反演变换，可以还原回 $\Gamma$。
    \item 反演中心$O$与无穷远点$\infty$互为反演点。
    \item 反演具有保角性，两曲线交角在反演前后保持不变，正交关系也保持不变。
\end{itemize}

在平面直角坐标系中，设反演中心为原点$O(0,0)$。对于任意一点$(x_0, y_0)$，其反演点坐标为$\left(\frac{R^{2}x_{0}}{x_{0}^{2}+y_{0}^{2}}, \frac{R^{2}y_{0}}{x_{0}^{2}+y_{0}^{2}}\right)$。根据反演的可逆性，对于一般二次曲线方程
$$Ax^{2}+Bxy+Cy^{2}+Dx+Ey+F=0$$
将其中的$x$替换为$\frac{R^{2}x}{x^{2}+y^{2}}$，将$y$替换为$\frac{R^{2}y}{x^{2}+y^{2}}$，即可得到反演像的代数方程：
$$R^{4}(Ax^{2}+Bxy+Cy^{2})+R^{2}(Dx+Ey)(x^{2}+y^{2})+F(x^{2}+y^{2})^{2}=0$$
其中 $x, y$ 不同时为 $0$。若 $F\ne0$，最高项为 $F(x^2+y^2)^2$；若 $F=0$，最高次数降至三次。直线和圆属于低次数情形，可单独处理。

例如，以原点$O$为反演中心，$R=2$为反演半径，对椭圆$\frac{x^{2}}{4}+y^{2}=1$进行反演。代入坐标变换公式并化简，可得反演像方程为$(x^{2}+y^{2})^{2}=4x^{2}+16y^{2}$。

\subsection{四次曲线的垂足曲线生成方式}
同一类四次曲线也可以从垂足轨迹写出。考虑椭圆 $\frac{x^2}{a^2}+\frac{y^2}{b^2}=1$，$O$ 是坐标原点，$l$ 是椭圆的动切线，$P$ 是点 $O$ 在直线 $l$ 上的垂足。设 $P(x_0, y_0)$，则切线 $l$ 的方程为 $y - y_0 = -\frac{x_0}{y_0}(x - x_0)$，整理得 $\frac{x_0}{x_0^2 + y_0^2}x + \frac{y_0}{x_0^2 + y_0^2}y = 1$。由相切条件可知 $a^2\left(\frac{x_0}{x_0^2 + y_0^2}\right)^2 + b^2\left(\frac{y_0}{x_0^2 + y_0^2}\right)^2 = 1$。整理得 $P$ 的轨迹方程为 $(x^2 + y^2)^2 = a^2x^2 + b^2y^2$。这正是椭圆 $\frac{x^2}{a^2}+\frac{y^2}{b^2}=1$ 在单位反演下的像。同理，对于双曲线 $xy = 1$，点 $O$ 到其切线的垂足轨迹方程为 $(x^2 + y^2)^2 = 4xy$，这就是双曲线 $xy = 1$ 以 $O$ 为中心、$\sqrt{2}$ 为半径时的反演像。

\section{基础圆锥曲线定理的反演命题}

利用反演变换的保角性以及点、线、圆之间的对应关系，可以把一些解析几何问题写成高次曲线的交点问题。

\begin{example}
已知曲线$\Gamma: (x^2 + y^2)^2 = 4x^2 + 16y^2$，$A(-2,0)$。过原点$O$和定点$E\left(-\frac{10}{3},0\right)$的动圆与$\Gamma$交于$B,C$两点，$\triangle AOB$与$\triangle AOC$的外接圆半径分别为$R_1,R_2$。证明：$(R_1^2 - 1)(R_2^2 - 1)$为定值。
\end{example}
\begin{solution}
设过点$O(0,0)$和$E\left(-\frac{10}{3},0\right)$的圆方程为
$$x^{2}+y^{2}+\frac{10}{3}x+Fy=0$$
联立曲线$\Gamma$的方程$(x^{2}+y^{2})^{2}=4x^{2}+16y^{2}$，消去二次项得$\left(-\frac{10}{3}x-Fy\right)^{2}=4x^{2}+16y^{2}$，整理得
$$64x^{2}+60Fxy+9(F^{2}-16)y^{2}=0$$
设$B(x_{1},y_{1})$，$C(x_{2},y_{2})$，且$y_{1},y_{2}\neq 0$。令$t_{1}=\frac{x_{1}}{y_{1}}$，$t_{2}=\frac{x_{2}}{y_{2}}$，则$t_{1},t_{2}$为方程$64t^{2}+60Ft+9(F^{2}-16)=0$的两根。由韦达定理可得
$$t_{1}+t_{2}=-\frac{15F}{16},\quad t_{1}t_{2}=\frac{9(F^{2}-16)}{64}$$
由于$\triangle ABO$的外接圆过原点$O$及定点$A(-2,0)$，可设其方程为$x^{2}+y^{2}+2x+E_{1}y=0$。将点$B$满足的圆方程$x_{1}^{2}+y_{1}^{2}=-\frac{10}{3}x_{1}-Fy_{1}$代入，得$-\frac{4}{3}x_{1}+(E_{1}-F)y_{1}=0$。由此解得$E_{1}=F+\frac{4}{3}t_{1}$。同理，$\triangle ACO$外接圆的参数满足$E_{2}=F+\frac{4}{3}t_{2}$。
两外接圆的半径平方分别为$R_{1}^{2}=1+\frac{E_{1}^{2}}{4}$与$R_{2}^{2}=1+\frac{E_{2}^{2}}{4}$。因此
$$(R_{1}^{2}-1)(R_{2}^{2}-1)=\frac{E_{1}^{2}E_{2}^{2}}{16}$$
计算$E_{1}E_{2}$的乘积：
$$E_{1}E_{2}=F^{2}+\frac{4}{3}F(t_{1}+t_{2})+\frac{16}{9}t_{1}t_{2}$$
代入韦达定理结论可得$E_{1}E_{2}=F^{2}-\frac{5}{4}F^{2}+\frac{1}{4}(F^{2}-16)=-4$。于是$(R_{1}^{2}-1)(R_{2}^{2}-1)=\frac{(-4)^{2}}{16}=1$，即为所求。
这也可理解为：在椭圆 $\frac{x^2}{4}+y^2=1$ 中，过点 $\left(-\frac{6}{5},0\right)$ 的弦，对左顶点张角为直角。
\end{solution}

\begin{example}
已知曲线$\Gamma: (x^2 + y^2)^2 = 4xy$。过坐标原点$O$的动圆$\odot Q$与$\Gamma$相切，$\odot Q$分别交$x$轴、$y$轴于$B, C$两点。求证：$\triangle OBC$的面积为定值。
\end{example}
\begin{solution}
设$\odot Q$方程为$x^2+y^2=Dx+Ey$。联立$\Gamma$方程，得$(Dx+Ey)^2=4xy$。整理得关于$\frac{x}{y}$的二次方程：
$$D^2\left(\frac{x}{y}\right)^2+(2DE-4)\left(\frac{x}{y}\right)+E^2=0$$
由相切条件，该方程的判别式必定为零：
$$\Delta=(2DE-4)^2-4D^2E^2=0$$
化简得$-16DE+16=0$，即$DE=1$。
圆$\odot Q$交坐标轴于$B(D,0)$及$C(0,E)$。$\triangle OBC$的面积为$\frac{1}{2}|D||E|=\frac{1}{2}$，即为所求。
\end{solution}

\section{进阶命题：蒙日圆与垂直弦性质的反演}

\begin{example}
已知曲线$\Gamma: (x^2 + y^2)^2 = 4x^2 + 16y^2$，$O$为坐标原点。有两个过原点$O$的圆$\odot C_1$与$\odot C_2$均与曲线$\Gamma$相切，且$\odot C_1$与$\odot C_2$在原点正交。设$\odot C_1$与$\odot C_2$除原点$O$外的交点为$P$。求证：线段$OP$的长度为定值。
\end{example}
\begin{solution}
设过原点的圆方程为$x^2 + y^2 = Dx + Ey$。将其与$\Gamma$的方程联立，考察相切条件，得$(Dx + Ey)^2 = 4x^2 + 16y^2$。整理得关于$x,y$的齐次二次方程：
$$(D^2 - 4)x^2 + 2DExy + (E^2 - 16)y^2 = 0$$
相切要求此二次齐次方程的判别式为零：
$$\Delta = (2DE)^2 - 4(D^2 - 4)(E^2 - 16) = 0$$
化简得$16D^2 + 4E^2 = 64$，即$4D^2 + E^2 = 16$。
设$\odot C_1, \odot C_2$的方程分别为$x^2 + y^2 = D_1x + E_1y$与$x^2 + y^2 = D_2x + E_2y$。两圆相切的条件为$4D_1^2 + E_1^2 = 16$且$4D_2^2 + E_2^2 = 16$。两圆在原点正交的条件为$D_1D_2 + E_1E_2 = 0$。
联立两圆方程，求交点$P(x,y)$。令$|OP|^2 = x^2 + y^2$。由两圆方程得$D_1x + E_1y = |OP|^2$且$D_2x + E_2y = |OP|^2$。运用克莱姆法则解得$x, y$，并代入$x^2 + y^2 = |OP|^2$中，化简得：
$$|OP|^2 = \frac{(D_1E_2 - D_2E_1)^2}{(D_1 - D_2)^2 + (E_1 - E_2)^2}$$
利用恒等式及正交条件$D_1D_2 + E_1E_2 = 0$，将分子和分母展开：
$$|OP|^2 = \frac{(D_1^2 + E_1^2)(D_2^2 + E_2^2)}{D_1^2 + E_1^2 + D_2^2 + E_2^2}$$
代入条件 $E_i^2=16-4D_i^2\ (i=1,2)$。记
\[
A=D_1^2+D_2^2,\qquad B=D_1^2D_2^2.
\]
由正交条件平方得
\[
B=E_1^2E_2^2=(16-4D_1^2)(16-4D_2^2),
\]
即
\[
15B=64A-256.
\]
又
\[
D_i^2+E_i^2=16-3D_i^2,
\]
所以
\[
\begin{aligned}
(D_1^2+E_1^2)(D_2^2+E_2^2)
&=(16-3D_1^2)(16-3D_2^2)\\
&=256-48A+9B\\
&=\frac{16}{5}(32-3A),
\end{aligned}
\]
而
\[
D_1^2+E_1^2+D_2^2+E_2^2=32-3A.
\]
代回得 $|OP|^2=\frac{16}{5}$。因此线段 $OP$ 的长为 $\frac{4\sqrt5}{5}$。
\end{solution}

\begin{example}
已知曲线$\Gamma: x^3 + xy^2 - y^2 = 0$（蔓叶线）。设$A, B$是$\Gamma$上异于原点$O$的两个动点，且满足$OA \perp OB$。求证：$\triangle OAB$的外接圆恒过一个异于原点的定点。
\end{example}
\begin{solution}
设直线$OA$的方程为$y = kx$，则$OB$的方程为$y = -\frac{1}{k}x$。代入曲线$\Gamma$求点$A, B$坐标。对点$A$，得$x^3 + k^2x^3 - k^2x^2 = 0$，解得$x_A = \frac{k^2}{1+k^2}$，$y_A = \frac{k^3}{1+k^2}$。同理，将$k$替换为$-\frac{1}{k}$得$x_B = \frac{1}{1+k^2}$，$y_B = -\frac{1}{k(1+k^2)}$。
设$\triangle OAB$外接圆方程为$x^2 + y^2 + Dx + Ey = 0$。代入点$A$坐标并除以$\frac{k^2}{1+k^2}$得$k^2 + D + Ek = 0$。代入点$B$坐标得$-\frac{1}{k^2} + D - \frac{E}{k} = 0$，即$Dk^2 - Ek = -1$。两式相加消去$E$得$D(1+k^2) = -(1+k^2)$，由此解得$D=-1$及$E=\frac{1-k^2}{k}$。外接圆方程为$x^2 + y^2 - x + \frac{1-k^2}{k}y = 0$。令含有参数的项$y=0$，得$x^2 - x = 0$。解得异于原点的定点为$(1, 0)$。
\end{solution}

\section{离心率、极点与极线的反演表示}

\begin{example}
已知曲线$\Gamma: (x^2+y^2+2x)^2 = 4(x^2+y^2)$（心脏线），$O$为坐标原点。过原点$O$作任意一条不与$y$轴重合的直线，交曲线$\Gamma$于$A, B$两点。求证：线段$AB$的长度为定值$4$。
\end{example}
\begin{solution}
设直线方程为$y = kx$。代入曲线$\Gamma$方程，得$(x^2 + k^2x^2 + 2x)^2 = 4(x^2 + k^2x^2)$。由于交点异于原点，提取公因式$x^2$并化简得$[x(1+k^2) + 2]^2 = 4(1+k^2)$。开平方后解出两交点的横坐标：
$$x_A = \frac{-2 + 2\sqrt{1+k^2}}{1+k^2}, \quad x_B = \frac{-2 - 2\sqrt{1+k^2}}{1+k^2}$$
线段长度$|AB| = \sqrt{1+k^2} |x_A - x_B| = \sqrt{1+k^2} \left| \frac{4\sqrt{1+k^2}}{1+k^2} \right| = 4$。
\end{solution}

极性变换与反演也能反映极点与极线中的调和共轭关系。

\begin{example}
已知曲线$\Gamma_1: (x^2+y^2)^2 - 2x(x^2+y^2) + 2y^2 = 0$以及圆$\Gamma_2: x^2+y^2 - x = 0$。过原点$O$作任意一条射线，交$\Gamma_1$于$A, B$两点，交$\Gamma_2$于$Q$点。求证：点$Q$恒为线段$AB$的中点。
\end{example}
\begin{solution}
设射线方程为 $y = kx$（$x > 0$）。将它代入圆 $\Gamma_2$ 的方程，得
\[
x(1+k^2)=1,
\]
所以
\[
x_Q=\frac{1}{1+k^2}.
\]
再代入曲线 $\Gamma_1$，约去非零因子后得关于横坐标的方程
\[
x^2(1+k^2)^2-2x(1+k^2)+2k^2=0.
\]
由韦达定理，
\[
x_A+x_B=\frac{2}{1+k^2}=2x_Q.
\]
由于 $A,B,Q$ 都在同一条射线上，故点 $Q$ 恰为线段 $AB$ 的中点。
\end{solution}

对于任意离心率$e > 0$，帕斯卡蜗牛线族$\Gamma_e: (x^2+y^2+ex)^2 = x^2+y^2$中过极点的弦长恒为常数。设截线$y = kx$，代入得$[x(1+k^2) + e]^2 = 1+k^2$。非原点两交点的横坐标之差为$\frac{2}{\sqrt{1+k^2}}$，再乘上直线方向因子$\sqrt{1+k^2}$，弦长便是 $2$，与离心率参数 $e$ 无关。

\section{射影几何定理的反演表示}

反演不仅适用于度量性质的推广，对于射影几何中的共点、共线等关系也能得到相应结论。

以帕斯卡定理为例：圆锥曲线上内接六边形对边延长线的三个交点共线。通过选定特定反演中心，双曲线反演生成双纽线$\Gamma: (x^2+y^2)^2 = 4xy$。原图形中的割线反演为过原点的圆，三点共线反演为四点共圆。由此得出结论：在$\Gamma$上取$6$个异于原点$O$的点，作$6$个过原点的圆，其对位交点$P, Q, R$必与点$O$共圆。

布里昂雄定理也可由极性变换得到；彭赛列闭合定理可用同一极性语言表述。这里不展开完整定理，只记录反演下需要用到的点线圆对应。

\section{隐蔽的反演命题：弗雷歇定理与中点轨迹}

以下题目都围绕反演、极性与弦中点轨迹之间的代数联系展开。

\begin{example}
已知曲线$\Gamma: 12x^3 + 12xy^2 + 3x^2 + 4y^2 = 0$。过原点$O$作两条互相垂直的直线，分别交$\Gamma$于异于原点的两点$A, B$。求证：$\triangle OAB$的外接圆恒过一定点。
\end{example}
\begin{solution}
设$OA$方程为$y = kx$，$OB$方程为$y = -\frac{1}{k}x$。代入$\Gamma$并降次，解得$x_A = -\frac{3+4k^2}{12(1+k^2)}$，$y_A = -\frac{k(3+4k^2)}{12(1+k^2)}$。同理得$x_B = -\frac{3k^2+4}{12(k^2+1)}$，$y_B = \frac{3k^2+4}{12k(k^2+1)}$。设外接圆方程为$x^2 + y^2 + Dx + Ey = 0$。代入$A, B$坐标并消元，化简可解得关于$D, E$的线性方程组。其中$D + Ek = \frac{3+4k^2}{12}$与$Dk^2 - Ek = \frac{3k^2+4}{12}$。两式相加消去$E$项，得$D(1+k^2) = \frac{7(1+k^2)}{12}$，即$D = \frac{7}{12}$。令含参项$y=0$，解得外接圆恒过的定点为$\left(-\frac{7}{12}, 0\right)$。
\end{solution}

\begin{example}
已知曲线$\Gamma: (x^2+y^2)^2 + 2x(x^2+y^2) - 2y^2 = 0$。过原点$O$的直线$l$交$\Gamma$于两点$A, B$，设中点为$M$。求$M$的轨迹。
\end{example}
\begin{solution}
设直线 $l$ 的方程为 $y = kx$。代入 $\Gamma$ 的方程，可得一元二次方程
\[
(1+k^2)^2x^2+2(1+k^2)x-2k^2=0.
\]
由韦达定理，中点 $M$ 的坐标为
\[
x_M=\frac{-1}{1+k^2},\qquad y_M=\frac{-k}{1+k^2}.
\]
消去参数 $k$，得
\[
x_M^2+y_M^2=\frac{1}{1+k^2}=-x_M.
\]
因此轨迹方程为
\[
x^2+y^2+x=0
\]
（除去原点）。这里也反映了反演前后点线关系的一种对应。
\end{solution}

\section{复平面上的反演变换}

在复平面$\mathbb{C}$中，取单位圆为反演基准，坐标$(x,y)$在复反演映射下可写成简单的代数形式：
$$w = \frac{1}{\bar{z}}$$
广义圆与直线的一般方程
\[
A z \bar{z} + \bar{B} z + B \bar{z} + C = 0
\]
在上述映射下化为
\[
C w \bar{w} + B \bar{w} + \bar{B} w + A = 0.
\]
可见二次项系数 $A$ 与常数项 $C$ 交换了位置，正好对应复平面反演的代数形式。

\begin{example}
已知非零复数$z_1, z_2$满足方程$z\bar{z} = z + \bar{z}$。设复数$z_3$满足$\frac{2}{z_3} = \frac{1}{z_1} + \frac{1}{z_2}$。关于$z_1, z_2, z_3$的论断如下：
A. 复数$z_3$必定满足方程$z\bar{z} = z + \bar{z}$ \\
B. $\frac{|z_1 - z_3|}{|z_1|} = \frac{|z_2 - z_3|}{|z_2|}$ \\
C. $\frac{1}{|z_1|} + \frac{1}{|z_2|} \ge \frac{2}{|z_3|}$ \\
D. $\frac{1}{|z_1|^2} + \frac{1}{|z_2|^2} \ge \frac{2}{|z_3|^2}$ \\
判断上述论断的正确性。
\end{example}
\begin{solution}
将原方程$z\bar{z} = z + \bar{z}$两边同除以$z\bar{z}$，得$\frac{1}{\bar{z}} + \frac{1}{z} = 1$。转化至反演平面中，令$w = \frac{1}{z}$，方程化为$w + \bar{w} = 1$，即$\text{Re}(w) = \frac{1}{2}$。该方程在$w$平面上表示一条垂直于实轴的直线。由调和平均条件$\frac{2}{z_3} = \frac{1}{z_1} + \frac{1}{z_2}$可知$w_3 = \frac{w_1 + w_2}{2}$，即$w_3$为线段$w_1 w_2$的中点。

对于论断A：由于$w_1, w_2$均在直线$\text{Re}(w) = \frac{1}{2}$上，其所在线段的中点$w_3$必满足$\text{Re}(w_3) = \frac{1}{2}$。映射回原复平面，$z_3$满足方程$z\bar{z} = z + \bar{z}$。A正确。

对于论断B：将距离公式转化至$w$平面。$\frac{|z_1 - z_3|}{|z_1|} = \frac{|w_3 - w_1|/|w_1 w_3|}{1/|w_1|} = \frac{|w_3 - w_1|}{|w_3|}$。同理等式右侧为$\frac{|w_3 - w_2|}{|w_3|}$。由中点性质可知$|w_3 - w_1| = |w_3 - w_2|$，故等式成立。B正确。

对于论断C：在$w$平面应用三角不等式。由于$w_1 + w_2 = 2w_3$，取模长得$|w_1| + |w_2| \ge |w_1 + w_2| = 2|w_3|$。将其还原为原坐标表达即为$\frac{1}{|z_1|} + \frac{1}{|z_2|} \ge \frac{2}{|z_3|}$。C正确。

对于论断D：由复数运算的平行四边形法则，$|w_1 - w_2|^2 + |w_1 + w_2|^2 = 2(|w_1|^2 + |w_2|^2)$。代入$w_1 + w_2 = 2w_3$并利用$|w_1 - w_2|^2 \ge 0$放缩可得$4|w_3|^2 \le 2(|w_1|^2 + |w_2|^2)$。还原后不等式成立。D正确。
\end{solution}

\begin{example}
已知曲线$\Gamma: (x^2+y^2)^2 = x^2-y^2$。设过坐标原点$O$的动圆$C$与$\Gamma$相切。若圆$C$分别交直线$y=x$与$y=-x$于异于原点的$M, N$两点，求证：$\triangle OMN$的面积为定值。
\end{example}
\begin{solution}
设圆$C$的方程为$x^2+y^2+Dx+Ey=0$。将该方程与直线$y=x$联立，可得交点纵横坐标满足$2x^2+(D+E)x=0$。若交点$M$异于原点，则$x_M=y_M=-\frac{D+E}{2}$。由此可得线段长度$|OM|=\frac{|D+E|}{\sqrt{2}}$。
同理，将其与直线$y=-x$联立，解得交点$N$满足$x_N=-\frac{D-E}{2}, y_N=\frac{D-E}{2}$，线段长度$|ON|=\frac{|D-E|}{\sqrt{2}}$。
因直线$y=x$与$y=-x$垂直，$\triangle OMN$面积表达为$S = \frac{1}{2}|OM||ON| = \frac{1}{4}|D^2-E^2|$。
由圆方程得 $x^2+y^2=-Dx-Ey$，代入 $\Gamma$ 得
\[
(-Dx-Ey)^2=x^2-y^2.
\]
整理为$(D^2-1)x^2 + 2DExy + (E^2+1)y^2 = 0$。
若相切则此式判别式等于零，即$(2DE)^2 - 4(D^2-1)(E^2+1) = 0$。化简可得方程约束$D^2 - E^2 = 1$。
将该约束代入面积公式中，得到$S = \frac{1}{4}$，该面积恒为定值。（本题几何背景为等轴双曲线切线交其渐近线所成三角形面积定值性质的反演）。
\end{solution}

\begin{example}
已知曲线$\Gamma: x^3+xy^2-y^2=0$。设有两个过原点$O$的动圆$C_1, C_2$均与$\Gamma$相切，且$C_1$与$C_2$在原点处正交。若$P$为两圆除原点外的交点，求证：点$P$恒在一个定圆上。
\end{example}
\begin{solution}
设切圆一般方程为$x^2+y^2+Dx+Ey=0$。代入$\Gamma$提取共因式后化为$Dx^2+Exy+y^2=0$。
相切要求判别式等于零，即满足$4D=E^2$。
设圆$C_1, C_2$的对应系数为$(D_1, E_1)$与$(D_2, E_2)$。两圆在原点正交等效于一次项系数内积$D_1D_2+E_1E_2=0$。代入切线约束条件后得
\[
\frac{E_1^2E_2^2}{16}+E_1E_2=E_1E_2\left(\frac{E_1E_2}{16}+1\right)=0.
\]
若 $E_1E_2=0$，则某个切圆满足 $E=0,D=0$，圆方程退化为 $x^2+y^2=0$，不符合“除原点外另有交点”的条件。因此只能取非零分支，得$E_1E_2 = -16$。
由方程作差联立解交点$P(x,y)$，得到直线约束$y = -\frac{D_1-D_2}{E_1-E_2}x = -\frac{E_1+E_2}{4}x$。
由 $E_1E_2=-16$ 和和式可得一元二次方程，进而有 $E_1^2 = -\frac{4y}{x}E_1 + 16$。
将$D_1=\frac{E_1^2}{4}$回代至$C_1$方程中并将二次形式带入替换，方程变为$x^2+y^2+\frac{1}{4}(-\frac{4y}{x}E_1 + 16)x+E_1y=0$。展开并合并同类项，含参项$E_1$相互抵消，恒等化简为$x^2+y^2+4x=0$。该定圆即为点$P$之轨迹。
\end{solution}

\begin{example}
已知曲线$\Gamma: 27x^2+36y^2+6x(x^2+y^2)-(x^2+y^2)^2=0$。设过坐标原点$O$的动圆$C$与$\Gamma$相切，且线段$OM$为圆$C$的一条直径。求证：点$M$的轨迹是一个圆。
\end{example}
\begin{solution}
设动圆$C$方程为$x^2+y^2-Dx-Ey=0$。因$OM$为其直径且过原点，端点$M$坐标恰为$(D, E)$。
代入曲线$\Gamma$可得关于$x,y$的二次齐次形式：$(27+6D-D^2)x^2 + (6E-2DE)xy + (36-E^2)y^2 = 0$。
结合相切条件判别式设为零，即$(6E-2DE)^2 - 4(27+6D-D^2)(36-E^2) = 0$。
展开判别式左端：
\[
(6E-2DE)^2 - 4(27+6D-D^2)(36-E^2)
=144(D^2-6D+E^2-27).
\]
因此有$36E^2 + 36D^2 - 216D - 972 = 0$。
对其同除以$36$，得
\[
D^2+E^2-6D-27=0.
\]
将坐标$(x_M, y_M)$替代$(D, E)$，可知点$M$的轨迹为圆$x^2+y^2-6x-27=0$。
\end{solution}

\begin{example}
定圆$C_1: x^2+y^2+4x=0$。对于$C_1$上任一异于原点$O$的点$P(x_0, y_0)$，作经过$O$与$P$且同曲线$\Gamma: x^3+xy^2-y^2=0$相切的两圆，设切点为$A, B$。求证：$\triangle OAB$的外接圆恒过一定点。
\end{example}
\begin{solution}
设经过$O, P$且与$\Gamma$相切的圆方程为$x^2+y^2+Dx+Ey=0$，其相切约束为$4D=E^2$。
由经过点$P(x_0, y_0)$，将已知定圆约束$x_0^2+y_0^2=-4x_0$带入切圆方程化简得$(D-4)x_0+Ey_0=0$。
代入$D$表达式构成二次方程$(E^2-16)x_0+4Ey_0=0$。其两根$E_1, E_2$由韦达定理给定$E_1E_2=-16$及$E_1+E_2=-\frac{4y_0}{x_0}$。
求解具体切点坐标。原齐次方程$Dx^2+Exy+y^2=0$给出切点所在的方向关系，写成$y=-\frac{E}{2}x$后，再与切圆方程联立，解得切点横纵坐标分别为$x = \frac{E^2}{E^2+4}, y = \frac{-2E}{E^2+4}$。
于是
\[
x_A^2+y_A^2
=\frac{E_1^4+4E_1^2}{(E_1^2+4)^2}
=\frac{E_1^2}{E_1^2+4}
=x_A,
\]
所以点$A$在圆$x^2+y^2-x=0$上。
同理，切点$B$坐标也满足该方程；原点$O$代入该方程也成立。因此$\triangle OAB$的外接圆方程为$x^2+y^2-x=0$，它经过定点$(1, 0)$。
\end{solution}



\begin{example}
曲线$\Gamma: x^3+xy^2-y^2=0$。圆$C$过原点$O$与定点$F(4,0)$，并且同$\Gamma$在另两点$A, B$相交。求证：直线$OA$与$OB$的斜率乘积为常数。
\end{example}
\begin{solution}
由动圆$C$经过点$O(0,0)$与$F(4,0)$两点，其圆心必位于直线$x=2$上。设圆$C$的一般方程为$x^2+y^2-4x+Ey=0$。
将圆$C$与曲线$\Gamma$联立以求相交弦属性。对于$\Gamma$方程可化为$x(x^2+y^2) = y^2$。
将圆$C$的方程代入，得到$x(4x-Ey) = y^2$，展开得$4x^2-Exy-y^2=0$。
两边同除以$x^2$转化为关于直线斜率$k=y/x$的方程，得到$k^2+Ek-4=0$。
该一元二次方程的两根$k_A, k_B$分别就是直线$OA$与$OB$的斜率。由韦达定理得
\[
k_Ak_B = -4.
\]
这个乘积与动参数$E$无关，因此斜率乘积恒为常数$-4$。
\end{solution}

复平面反演也可用于处理经典几何命题。以托勒密定理为例，若把反演中心取在圆内接四边形的一个顶点，原圆会化为一条直线，命题便可转化为共线三点的距离关系
\[
|B'D'|=|B'C'|+|C'D'|.
\]
再代入距离公式
\[
|w_X-w_Y|=\frac{|z_X-z_Y|}{|z_X||z_Y|}
\]
若反演中心对应原四边形顶点 $A$，则
\[
\frac{BD}{AB\cdot AD}
=\frac{BC}{AB\cdot AC}+\frac{CD}{AC\cdot AD}.
\]
两边同乘 $AB\cdot AC\cdot AD$，得到
\[
AC\cdot BD=AD\cdot BC+AB\cdot CD,
\]
即托勒密等式。

\begin{example}
设曲线 $\Gamma: x^3 + xy^2 - 3x^2 + 2y^2 = 0$。动圆 $C$ 过原点 $O(0,0)$ 与定点 $F(4,0)$，并且与曲线 $\Gamma$ 在另两点 $A, B$ 相交。求证：直线 $OA$ 与 $OB$ 的斜率乘积为常数。
\end{example}
\begin{solution}
动圆 $C$ 经过 $O(0,0)$ 和 $F(4,0)$，故圆心在直线 $x=2$ 上。设
\[
C:\ x^2+y^2-4x+Ey=0,
\]
即 $x^2+y^2=4x-Ey$。又
\[
\Gamma:\ x(x^2+y^2)-3x^2+2y^2=0.
\]
代入圆方程得
\[
x(4x-Ey)-3x^2+2y^2=0,
\]
即
\[
x^2-Exy+2y^2=0.
\]
除以 $x^2$，令 $k=\frac yx$，得到
\[
2k^2-Ek+1=0.
\]
该方程的两根就是 $OA,OB$ 的斜率 $k_A,k_B$，故
\[
k_Ak_B=\frac12.
\]
\end{solution}

\begin{example}
设另一曲线 $\Gamma: x^3 + xy^2 - x^2 - 2y^2 = 0$。动圆 $C$ 过原点 $O(0,0)$ 与定点 $P(0,2)$，并且与曲线 $\Gamma$ 相交于另外两点 $A, B$。求证：直线 $OA$ 与 $OB$ 的斜率之和为定值。
\end{example}
\begin{solution}
动圆 $C$ 经过 $O(0,0)$ 和 $P(0,2)$，故圆心在直线 $y=1$ 上。设
\[
C:\ x^2+y^2+Dx-2y=0,
\]
即 $x^2+y^2=2y-Dx$。又
\[
\Gamma:\ x(x^2+y^2)-x^2-2y^2=0.
\]
代入圆方程得
\[
x(2y-Dx)-x^2-2y^2=0,
\]
即
\[
-2y^2+2xy-(D+1)x^2=0.
\]
除以 $-x^2$，令 $k=\frac yx$，得到
\[
2k^2-2k+(D+1)=0.
\]
该方程的两根就是 $OA,OB$ 的斜率 $k_A,k_B$，故
\[
k_A+k_B=1.
\]
\end{solution}

\begin{example}
    已知曲线 $\Gamma: x^3 + xy^2 + 2x^2 + xy + y^2 = 0$。动圆 $C$ 过原点 $O(0,0)$ 与定点 $P(-1, 1)$，并且与曲线 $\Gamma$ 在另外两点 $A, B$ 相交。求证：直线 $OA$ 与 $OB$ 的斜率满足关系式 $k_A k_B + k_A + k_B$ 为常数，并求出该常数的值。
\end{example}
\begin{solution}
设经过原点 $O(0,0)$ 的动圆 $C$ 为
\[
x^2+y^2+Dx+Ey=0.
\]
动圆经过 $P(-1,1)$，代入得
\[
(-1)^2+1^2-D+E=0,
\]
即
\[
E-D=-2.
\]
于是
\[
x^2+y^2=-Dx-Ey.
\]
曲线 $\Gamma$ 可写成
\[
x(x^2+y^2)+2x^2+xy+y^2=0.
\]
代入动圆的二次形式，得
\[
x(-Dx-Ey)+2x^2+xy+y^2=0,
\]
整理为
\[
(2-D)x^2+(1-E)xy+y^2=0.
\]
除以 $x^2$，令 $k=\frac yx$，得到
\[
k^2+(1-E)k+(2-D)=0.
\]
该方程的两个根就是 $OA,OB$ 的斜率 $k_A,k_B$。由韦达定理，
\[
k_A+k_B=E-1,\qquad k_Ak_B=2-D.
\]
因此
\[
k_Ak_B+k_A+k_B=(2-D)+(E-1)=1+(E-D)=-1.
\]
故所求常数为 $-1$。
\end{solution}

\begin{example}
    已知曲线 $\Gamma: x^3 + xy^2 - y^2 = 0$。动直线 $L$ 不经过原点 $O(0,0)$，且与曲线 $\Gamma$ 交于 $A, B, C$ 三点。记直线 $OA, OB, OC$ 的斜率分别为 $k_A, k_B, k_C$。若这三条直线的斜率满足条件
    \[
    k_A + k_B + k_C + k_A k_B k_C = 2,
    \]
    求证：动直线 $L$ 恒过一个定点，并求出该定点坐标。
\end{example}
\begin{solution}
设不过原点的动直线 $L$ 的方程为
\[
ux+vy=1.
\]
曲线 $\Gamma$ 可写成
\[
x(x^2+y^2)=y^2.
\]
用直线方程代换右边的常数 $1$，得
\[
x(x^2+y^2)=y^2(ux+vy),
\]
即
\[
x^3+(1-u)xy^2-vy^3=0.
\]

交点 $A,B,C$ 都不在原点，且可设 $x\neq 0$。两边同除以 $x^3$，再令
\[
k=\frac{y}{x},
\]
得到
\[
1+(1-u)k^2-vk^3=0,
\]
即
\[
vk^3+(u-1)k^2-1=0.
\]
这个方程的三个根正是直线 $OA,OB,OC$ 的斜率 $k_A,k_B,k_C$。由韦达定理，
\[
k_A+k_B+k_C=\frac{1-u}{v},\qquad k_Ak_Bk_C=\frac{1}{v}.
\]
由题设条件
\[
k_A+k_B+k_C+k_Ak_Bk_C=2
\]
可得
\[
\frac{1-u}{v}+\frac{1}{v}=2,
\]
即
\[
\frac{2-u}{v}=2 \iff u+2v=2.
\]
因此
\[
u\cdot \frac12+v\cdot 1=1.
\]
与直线方程 $ux+vy=1$ 对比可知，动直线 $L$ 恒过定点
\[
\left(\frac12,1\right).
\]
\end{solution}
上述结论也可从反演角度理解。原图形中，抛物线 $y^2=x$ 上有一族过原点的动圆；若这些圆除了原点外还与抛物线交于另外三点，并且恒过定点 $(2,1)$，那么以原点为反演中心后，抛物线会变成三次曲线
\[
\Gamma:x^3+xy^2-y^2=0,
\]
过原点的圆会变成不过原点的直线，而点 $(2,1)$ 的反演点正是
\[
\left(\frac12,1\right).
\]
所以上面的定点结论正好可以理解为原结论的反演像。


\chapter{杂题}
\begin{example}[浙江风味圆锥曲线]
椭圆$C_1:\frac{x^2}4+\frac{y^2}3=1$,抛物线$C_2:y^2=2px(p>0).$斜率为正的直线$l$与椭圆$C_1$和抛物线$C_2$均恰有一个公共点，分别为点$A,B$,椭圆$C_1$交抛物线$C_2$ 于点 $C,D.$\\(1)若直线$CD$经过椭圆$C_1$的右焦点$F$,求$p$的值；\\(2)设直线$CD$交直线$l$于点$E$,求$\triangle AOE$面积的最小值.
\end{example}
\begin{solution}
已知椭圆 $C_1:\frac{x^2}{4}+\frac{y^2}{3}=1$ 的右焦点为 $F(c, 0)$，其中 $c = \sqrt{4-3} = 1$，即 $F(1,0)$。联立 $C_1$ 与 $C_2:y^2=2px$ 时，方程只含 $y^2$，因此两个交点 $C,D$ 的横坐标相同、纵坐标相反，直线 $CD$ 垂直于 $x$ 轴。若直线 $CD$ 经过右焦点 $F(1,0)$，则点 $C,D$ 的横坐标必定为 $x=1$。将 $x=1$ 代入椭圆 $C_1$ 的方程中，解得 $y^2 = \frac{9}{4}$。由于点 $C, D$ 亦在抛物线 $C_2$ 上，将 $x=1, y^2=\frac{9}{4}$ 代入抛物线方程 $y^2=2px$ 中，得 $\frac{9}{4} = 2p$，解得 $p = \frac{9}{8}$。

求 $\triangle AOE$ 面积的最小值。作伸缩变换
\[
x = 2X,\qquad y = \sqrt{3}Y.
\]
变换后，椭圆 $C_1$ 化为单位圆 $C_1': X^2+Y^2=1$，抛物线 $C_2$ 化为
\[
3Y^2=4pX.
\]
记 $p' = \frac{2p}{3}$，则抛物线方程可写成
\[
C_2': Y^2=2p'X.
\]
直线 $l$ 仍是 $C_1'$ 和 $C_2'$ 的公切线。设它交 $X$ 轴于点 $T(t,0)$。由伸缩变换可知
\[
S_{\triangle AOE}=2\sqrt{3}\,S_{\triangle A'O'E'}.
\]

设二次曲线方程为 $S(x,y)=0$，曲线外一点为 $P_0(x_0,y_0)$。记 $P_0$ 代入曲线方程所得的值为 $S_{11}$，对应极线方程为 $T(x,y)=0$，则从 $P_0$ 向该二次曲线所作两条切线的联合方程为
\[
T^2=S\cdot S_{11}.
\]
该式由切线判别式得到：取过 $P_0$ 的动直线与二次曲线相交，把交点参数代入 $S=0$ 后得到一个二次方程；相切时判别式为 $0$，整理即为上式。

由上述结论，设点 $T(t, 0)$，由于 $l$ 为公切线，从 $T$ 点引出的两条切线满足如下方程：
对于单位圆 $C_1'$，代入公式可得方程为 $(tX-1)^2 = (t^2-1)(X^2+Y^2-1)$。
对于抛物线 $C_2'$，其极线为 $-p'(X+t)=0$，代入点 $T$ 得 $S_{11} = -2p't$，其切线对方程为 $(-p'(X+t))^2 = (-2p't)(Y^2-2p'X)$，化简得 $p'(X+t)^2 = -2t(Y^2-2p'X)$。
因直线 $l$ 是公切线且 $T$ 位于 $X$ 轴上，从 $T$ 引出的两条切线关于 $X$ 轴成对出现：一条是公切线 $l$，另一条是它关于 $X$ 轴的对称线。因此圆与抛物线的切线对表示同一对直线。

把单位圆的切线对展开为
\[
X^2-(t^2-1)Y^2-2tX+t^2=0,
\]
其 $X^2,Y^2$ 系数比为 $\frac{1}{1-t^2}$。抛物线的切线对展开为
\[
p'X^2+2p'tX+p't^2+2tY^2-4p'tX=0,
\]
其 $X^2,Y^2$ 系数比为 $\frac{p'}{2t}$。两对直线相同，二次项比例相同，故
\[
\frac{1}{1-t^2}=\frac{p'}{2t},
\]
解得 $p' = \frac{2t}{1-t^2}$。

直线 $CD$ 为交点弦，联立圆 $X^2+Y^2=1$ 与抛物线 $Y^2=2p'X$，消去 $Y$ 可得交点横坐标 $X_0$ 满足 $X_0^2 + 2p'X_0 - 1 = 0$，解得 $p' = \frac{1-X_0^2}{2X_0}$。将两个关于 $p'$ 的表达式联立，得 $\frac{1-X_0^2}{2X_0} = \frac{2t}{1-t^2}$，交叉相乘并展开重组得 $(X_0+t)^2 = (1-X_0t)^2$。由几何关系可知交点 $T$ 在抛物线左侧，即 $t<0$，且交点 $X_0 \in (0,1)$。若取
\[
X_0+t=1-X_0t,
\]
则 $t=\frac{1-X_0}{1+X_0}>0$，不合。故取
\[
X_0+t=-(1-X_0t),
\]
于是
\[
t=\frac{X_0+1}{X_0-1}.
\]

在变换后的图形中，设圆心为 $O'(0,0)$，$A'$ 为切点，则 $O'A' \perp A'E'$ 且半径 $|O'A'|=1$。在直角三角形 $\triangle O'A'E'$ 中，
\[
(S_{\triangle A'O'E'})^2
=\frac14 |O'A'|^2|A'E'|^2
=\frac14(X_0^2+Y_0^2-1).
\]
因交点 $E'(X_0,Y_0)$ 位于从 $T$ 引出的圆切线对上，代入
\[
(tX-1)^2=(t^2-1)(X^2+Y^2-1)
\]
得
\[
X_0^2+Y_0^2-1=\frac{(tX_0-1)^2}{t^2-1}.
\]
再代入 $t=\frac{X_0+1}{X_0-1}$，有
\[
tX_0-1=\frac{X_0^2+1}{X_0-1},\qquad
t^2-1=\frac{4X_0}{(X_0-1)^2}.
\]
代回面积公式，得到
\[
(S_{\triangle A'O'E'})^2=\frac{(X_0^2+1)^2}{16X_0}.
\]

设函数 $f(X_0) = \frac{(X_0^2+1)^2}{16X_0}$，其中 $X_0 \in (0, 1)$。对其求导得 $f'(X_0) = \frac{16(X_0^2+1)(3X_0^2-1)}{256X_0^2}$。令 $f'(X_0) = 0$，在区间 $(0,1)$ 内解得唯一的驻点 $X_0 = \frac{\sqrt{3}}{3}$。当 $X_0 \in (0, \frac{\sqrt{3}}{3})$ 时，$f'(X_0) < 0$，$f(X_0)$ 单调递减；当 $X_0 \in (\frac{\sqrt{3}}{3}, 1)$ 时，$f'(X_0) > 0$，$f(X_0)$ 单调递增。由此可知 $f(X_0)$ 在 $X_0 = \frac{\sqrt{3}}{3}$ 处取得最小值，计算得 $f\left(\frac{\sqrt{3}}{3}\right) = \frac{\sqrt{3}}{9}$，即 $S_{\triangle A'O'E'}$ 的最小值为 $\frac{\sqrt[4]{3}}{3}$。

根据仿射变换还原面积的几何倍数关系，可得原三角形面积的最小值为 $S_{min} = 2\sqrt{3} \times S_{\triangle A'O'E'} = 2\sqrt{3} \times \frac{\sqrt[4]{3}}{3} = \frac{2}{\sqrt[4]{3}}$。故 $\triangle AOE$ 面积的最小值为 $\frac{2}{\sqrt[4]{3}}$。
\end{solution}


\chapter{基本空间与元素}

中学里很少成体系讲射影几何，但它和解析几何、圆锥曲线、透视图法，以及许多“定点定线”问题联系很紧。它把平行、无穷远方向、二次曲线、极点极线、调和点列等现象放在一起讨论。

本章使用以下记号处理无穷远元素、齐次坐标、交比、极点极线与对偶关系。

在圆锥曲线问题中，极点极线常用于处理“过定点”“在定直线”“定值”等结论；它也会与交比、调和共轭、完全四边形、对偶原理、二次型和射影变换配合使用。

以下内容按三类对象组织：

\begin{enumerate}
    \item \textbf{直观几何视角}：从平行线、透视、无穷远点与投影出发，先建立对射影空间的几何理解；
    \item \textbf{代数坐标视角}：用齐次坐标、矩阵、线性变换与二次型，把射影几何写成代数形式；
    \item \textbf{圆锥曲线应用视角}：把极点极线、配极原则、自极三角形、焦点准线等问题放在同一思路下理解。
\end{enumerate}

以下内容限于圆锥曲线中常用的定义、公式与基本应用。

\begin{conclusion}
射影几何里最常用的几件事是：无穷远元素、点线对偶、交比、二次曲线。\\
把这些放在一起用，许多分散出现的结论就能更顺地串起来。
\end{conclusion}

\section{为什么要从仿射空间走向射影空间}
普通解析几何研究的是仿射空间。例如平面直角坐标系对应的就是二维仿射空间 $\A^2(\R)$。在这个空间里，两点确定一条直线，直线可以写成一次方程，圆锥曲线可以写成二次方程。但有一个最直接的问题：平行线没有交点。

射影几何的做法就是给每个方向补一个无穷远点，让同方向的平行直线都经过它。这样，原来的仿射平面就成了射影平面，所有无穷远点合在一起就是无穷远直线。

于是，原来的仿射平面就成了一个射影平面。此时：

\begin{itemize}
    \item 任意两条不同直线都相交于唯一一点；
    \item 平行不再是特殊情形，而只是“交点落在无穷远直线上”的普通情形；
    \item 椭圆、抛物线、双曲线在无穷远处的区别，变成了它们与无穷远直线交点个数的区别。
\end{itemize}

所以，射影空间可以看作在仿射空间中补入无穷远结构得到的几何，很多命题都能放在同一框架下讨论。

\begin{definition}
设 $\F$ 是一个域。考虑向量空间 $\F^{n+1}$ 中所有非零向量的集合
\[
\F^{n+1}\setminus\{ 0\}.
\]
在其上定义等价关系
\[
(x_0,x_1,\dots,x_n)\sim (y_0,y_1,\dots,y_n)
\quad \Longleftrightarrow \quad
\exists \lambda\in \F^\times,\ (y_0,\dots,y_n)=\lambda(x_0,\dots,x_n).
\]
这些等价类所构成的集合称为 \textbf{$n$维射影空间}，记作
\[
\PP^n(\F).
\]
\end{definition}

这一定义的含义是：射影空间中的“点”不是某个具体向量，而是\textbf{一条过原点的直线}。因为同一直线上的所有非零向量只差一个非零倍数，在射影空间里将它们视为同一点。

\section{射影空间、射影平面与射影直线}

\subsection{实射影空间与复射影空间}
当底层域 $\F=\R$ 时，得到的是\textbf{实射影空间}：
\[
\RP^n=\PP^n(\R).
\]
当底层域 $\F=\C$ 时，得到的是\textbf{复射影空间}：
\[
\CP^n=\PP^n(\C).
\]

对大多数解析几何与中学背景下的问题，实射影空间已经够用。若想把一些“在实数域下没有交点”的情形也一并讨论，常会把底层域换成复数域；例如两条平面二次曲线的交点总数在复射影平面里有一个固定的计数规律（把重根也按重数计）。

\subsection{射影平面}
\begin{definition}
二维射影空间 $\PP^2(\F)$ 称为 \textbf{射影平面}。特别地，$\RP^2$ 称为实射影平面，$\CP^2$ 称为复射影平面。
\end{definition}

射影平面可以写成：

\[
\text{射影平面}=\text{仿射平面}+\text{一条无穷远直线}.
\]

在仿射平面里，每个方向补一个无穷远点；所有方向对应的无穷远点合在一起，就得到无穷远直线。原来的仿射平面也就自然嵌入了射影平面。

\subsection{射影直线}
\begin{definition}
一维射影空间 $\PP^1(\F)$ 称为 \textbf{射影直线}。
\end{definition}

射影直线就是“普通直线再加上一个无穷远点”。对于实数域，可以写作
\[
\RP^1 \cong \R\cup\{\infty\}.
\]
这只是一个便于理解的写法：普通数轴向左右延伸到无穷远后，在射影直线里把两端接到同一个无穷远点上。

因此，射影直线可视为在普通数轴上补入一个无穷远点后的对象。点列、交比、调和点列和一维射影变换都在这条射影直线上表述。

\section{点、直线、无穷远点与无穷远直线}

\subsection{点：射影空间中的基本元素}
\begin{definition}
射影空间中的\textbf{点}是向量空间中一条过原点的一维线性子空间，常用一组不全为零的比例坐标
\[
(x_0:x_1:\cdots:x_n)
\]
来表示。
\end{definition}

这里的冒号不是普通有序数组的符号，而是在强调：这些坐标只在比例意义下确定。例如
\[
(1:2:3)=(2:4:6)=(-1:-2:-3).
\]
它们都表示同一个射影点。

射影坐标的基本思想是：\textbf{点不是由一个固定坐标表示，而是由一整类成比例的坐标共同表示。}

\subsection{齐次坐标}
\begin{definition}
设 $P\in \PP^n(\F)$。若 $P$ 由非零向量 $(x_0,x_1,\dots,x_n)$ 所对应的一维子空间表示，则称
\[
(x_0:x_1:\cdots:x_n)
\]
为点 $P$ 的 \textbf{齐次坐标}。
\end{definition}

齐次坐标的一个好处在于：有限点与无穷远点都可以用同样的方式表示。以实射影平面 $\RP^2$ 为例：

\begin{itemize}
    \item 若点的第三坐标 $z\neq 0$，则可化成 $(x:y:1)$ 的形式，对应仿射平面中的有限点 $(x,y)$；
    \item 若点满足 $z=0$，则它位于无穷远直线上，是一个无穷远点。
\end{itemize}

因此，用齐次坐标后，不必再把“无穷远点”当成例外去另立规则；它们只是射影点里的一类。

\subsection{直线：射影平面中的一维子空间}
在射影平面中，一条直线可以由一个齐次一次方程来表示。

\begin{definition}
在射影平面 $\PP^2(\F)$ 中，一个不全为零的齐次线性方程
\[
ax+by+cz=0
\]
所确定的点集称为一条 \textbf{直线}。
\end{definition}

在线性代数语言里，这条直线对应向量空间中的一个二维子空间；在坐标语言里，它是所有满足齐次一次方程的比例类点所组成的集合。

有一个基本事实：在射影平面里，\textbf{点与直线在代数形式上是对称的}。点写成比例类 $(x:y:z)$，直线写成系数比例类 $[a:b:c]$，二者之间的关联关系就是
\[
ax+by+cz=0.
\]
这种对称性正是对偶原理的根源。

\subsection{无穷远点}
\begin{definition}
仿射空间中同一方向的一族平行直线，在射影化之后所共有的交点称为该方向的 \textbf{无穷远点}。
\end{definition}

例如，在仿射平面中，所有斜率为 $k$ 的直线
\[
y=kx+b
\]
两两平行。在射影平面中，它们会交于同一个无穷远点
\[
(1:k:0).
\]
而所有形如 $x=c$ 的竖直直线，也会交于无穷远点
\[
(0:1:0).
\]

因此无穷远点并不是唯一的，而是“每个方向一个”。方向不同，无穷远点也不同。

\subsection{无穷远直线}
\begin{definition}
射影平面中所有无穷远点所组成的直线称为 \textbf{无穷远直线}，通常记作 $\ell_\infty$。
\end{definition}

在标准齐次坐标下，无穷远直线通常写作
\[
z=0.
\]

于是，仿射平面恰好可以被看作射影平面去掉无穷远直线以后剩下的部分：
\[
\A^2(\F)\cong \PP^2(\F)\setminus \ell_\infty.
\]

这条公式说明：仿射几何并不是和射影几何完全无关的另一套体系，而只是射影几何中“去掉无穷远直线”后剩下的部分。

\begin{property}
在射影平面中，任意两条不同直线都相交于唯一一点；其中，仿射平面中的平行线在射影平面中的交点恰好落在无穷远直线上。
\end{property}

\section{射影坐标与仿射坐标的关系}

\subsection{从仿射坐标到齐次坐标}
在二维情形下，仿射平面中的一个点 $(u,v)$ 可以嵌入射影平面为
\[
(u,v)\longmapsto (u:v:1).
\]
也可以写作 $(1:u:v)$，这取决于采用哪一种约定。这里采用
\[
(x,y)\mapsto (x:y:1)
\]
这一写法。

这种嵌入说明：仿射点只是射影点中第三坐标不为零的一类特殊点。而一旦第三坐标为零，就进入无穷远部分。

\subsection{齐次化}
很多仿射方程在射影几何中要改写成齐次方程。例如：

\begin{itemize}
    \item 仿射直线方程
    \[
    ax+by+c=0
    \]
    的齐次形式为
    \[
    ax+by+cz=0.
    \]
    \item 仿射二次曲线方程
    \[
    Ax^2+Bxy+Cy^2+Dx+Ey+F=0
    \]
    的齐次形式为
    \[
    Ax^2+Bxy+Cy^2+Dxz+Eyz+Fz^2=0.
    \]
\end{itemize}

这样就把无穷远处的情况也写进同一个方程里。于是，椭圆、抛物线、双曲线与无穷远直线的交点情况就能直接区分三类曲线。

\subsection{射影坐标}
\begin{definition}
在射影空间中，选定一组参考点或参考标架后所建立的比例坐标系统，称为 \textbf{射影坐标}。
\end{definition}

射影坐标与仿射坐标的区别在于：它不强调固定原点和长度单位，而强调按比例看同一方向，所以更适合处理射影几何里的比例关系。

从一维情形看，射影直线上只要固定三个互不重合的基准点，就可以把其他点参数化；从二维情形看，射影平面只要选定四个一般位置点，就能建立对应的射影坐标系。

\vspace{0.5em}
\begin{remark}
“射影坐标”和“齐次坐标”关系很近，但侧重方向不同。齐次坐标强调的是用比例类表示点，射影坐标强调的是在射影空间中建立坐标系统。
\end{remark}

\chapter{变换与映射}

\section{射影变换的基本思想}
欧氏几何主要研究刚体运动，所以长度、角度都保持不变；在仿射几何中允许的变换更大，于是平行和某些比值会保留下来；在射影几何中再放宽到中心投影等变换后，长度和角度一般不再保持，但共线、共点、交比等性质仍能保留。

射影变换通常从透视得到。把一个平面上的图形投到另一个平面上时，长度和角度一般都会变，但直线仍然映成直线。

\begin{definition}
设 $T:\F^{n+1}\to\F^{n+1}$ 是一个非奇异线性变换，则它在射影空间 $\PP^n(\F)$ 上诱导一个映射
\[
[x]\longmapsto [Tx].
\]
由此得到的点到点的一一对应称为 \textbf{射影变换} 或 \textbf{射影映射}。
\end{definition}

这里 $[x]$ 表示向量 $x$ 所对应的比例类。由于如果 $x$ 与 $\lambda x$ 表示同一点，则 $Tx$ 与 $\lambda Tx$ 仍表示同一点，所以这个定义是良定的。

\begin{property}
射影变换具有如下基本性质：
\begin{enumerate}
    \item 点映点，直线映直线；
    \item 共线点仍共线；
    \item 共点线仍共点；
    \item 交比保持不变；
    \item 二次曲线映为二次曲线。
\end{enumerate}
\end{property}

这些性质说明：射影变换并不是任意扭曲图形，它虽然通常不保持长度和角度，但会保留共线、共点、交比等基本关系。

\section{透视对应}
\begin{definition}
若两个点列、或两个线束、或一个点列与一个线束之间，通过某个固定中心或某条固定轴建立起对应关系，则称这种对应为 \textbf{透视对应}。
\end{definition}

比如，取平面上一点 $O$，用过 $O$ 的一束直线去截两条不同的直线 $l_1,l_2$。这样，$l_1$ 上的每个点都会通过与 $O$ 的连线对应到 $l_2$ 上的一个点，于是两个点列之间就形成了透视对应。

透视对应是射影变换的一个常见来源。许多一维结论（例如交比保持不变、调和点列保持不变、点列与线束的对应）都可以用透视来理解。

\begin{conclusion}
因此，一维射影几何里的点列与线束可以通过“投影—截影”互相转化，而射影不变量正是在这类转化中保持不变的量。
\end{conclusion}

\section{直射变换与对射变换}

\subsection{直射变换}
\begin{definition}
将点映为点、将直线映为直线，并保持关联关系的射影变换，称为 \textbf{直射变换}。
\end{definition}

通常所说的“射影变换”，默认就是这种点到点、线到线的变换。平面图形经过投影后，常得到这类变换。

\subsection{对射变换}
\begin{definition}
若一个映射把点映为直线，把直线映为点，并保持“点在直线上”这一关联关系的对偶形式，则称其为 \textbf{对射变换}。
\end{definition}

对射变换在此作为点线对应的代数语言使用；极点极线本身就是一种“点 $\leftrightarrow$ 直线”的对应。

直射变换处理的是点到点、线到线的对应；对射变换处理的则是点与直线之间的对应。射影几何中的对偶原理，正建立在这种点线互换上。

\section{对合}
\begin{definition}
一个变换 $f$ 若满足
\[
f\circ f=\mathrm{id},
\]
则称 $f$ 为一个 \textbf{对合}。
\end{definition}

对合的意思是：做一次变换得到某个像，再做一次同样的变换，就回到原对象。在线性代数中，关于某条直线的对称、关于某个点的中心对称都是对合；在射影几何中，一维射影直线上的许多重要变换也是对合。

对合在射影几何里经常和“成对出现”的对应一起出现。例如，在一条直线上，把每个点对应到某个共轭点；如果这个对应是互逆的，那它就是一个对合。极点极线、二次曲线上的共轭关系、完全四边形诱导出的调和对应，都属于这类情形。

\chapter{不变量与度量}

\section{什么是射影不变量}
几何学从根本上说，就是在某类允许的变换下研究保持不变的对象。不同几何学所研究的不变量不同：

\begin{center}
\begin{tabular}{>{\centering\arraybackslash}m{3cm}>{\centering\arraybackslash}m{4cm}>{\centering\arraybackslash}m{5cm}}
\hline
几何类型 & 允许变换 & 主要不变量 \\
\hline
欧氏几何 & 刚体运动 & 长度、角度 \\
相似几何 & 相似变换 & 角度、长度比 \\
仿射几何 & 仿射变换 & 共线、平行、单比 \\
射影几何 & 射影变换 & 共线、共点、交比 \\
\hline
\end{tabular}
\end{center}

因此，射影几何中主要的数值型不变量不是长度，也不是角度，而是\textbf{交比}。而像“共线”“共点”“点在线上”“线过某点”这样的关联性质，也都属于射影不变量。

\section{交比（复比）}
\begin{definition}
设直线上有四个有序点 $A,B,C,D$，则它们的 \textbf{交比}（或复比）定义为
\[
(A,B;C,D)=\frac{AC/BC}{AD/BD},
\]
其中线段采用有向线段的意义。
\end{definition}

若用仿射坐标表示，设四点坐标分别为 $a,b,c,d$，则
\[
(A,B;C,D)=\frac{(c-a)(d-b)}{(c-b)(d-a)}.
\]

交比并非只为当前计算设置，而是一维射影几何中的基本数值不变量。在射影变换之下，四个点的交比保持不变。

\begin{theorem}
在一维射影变换下，四个点的交比保持不变。
\end{theorem}

\begin{myproof}
设射影直线上的变换为
\[
x\longmapsto \frac{ax+b}{cx+d},\qquad ad-bc\neq 0.
\]
把四个点的坐标代入交比公式化简，便得到变换前后的交比相同。故交比是射影变换下的不变量。
\end{myproof}

因此在射影变换下，长度和角度一般不保，但交比保持不变，所以交比是最常用的数值不变量。

\section{调和共轭与调和点列}
\begin{definition}
若四个共线点 $A,B,C,D$ 的交比满足
\[
(A,B;C,D)=-1,
\]
则称这四点构成 \textbf{调和点列}，称 $C,D$ 关于 $A,B$ 互为 \textbf{调和共轭}。
\end{definition}

调和点列是交比的一个重要特例。原因在于：

\begin{enumerate}
    \item 它在射影变换下保持不变；
    \item 它会自然出现在完全四边形、极点极线和二次曲线的各种关系中；
    \item 它和“中点”“调和平均”“角平分线”等许多初等几何现象有关，常用来把这些结论写成同一种形式。
\end{enumerate}

若用有向线段写，调和点列还可改写为
\[
\frac{AB}{BC}=\frac{AD}{DC}.
\]
还可得到
\[
\frac{1}{AB}+\frac{1}{AD}=\frac{2}{AC},
\]
这也是“调和”一词的来源之一。

\vspace{0.5em}
\begin{property}
调和点列在透视对应下保持不变。若一点列是调和点列，则把它从某中心投影到另一条直线后，得到的对应四点仍构成调和点列。
\end{property}

这条性质是后续调和线束、完全四边形调和性质、极点极线几何定义以及配极原则的基础。

\section{射影度量}
射影几何里一般不把长度和角度当作基本量，因为它们在射影变换下通常不保持。不过，如果在射影空间里额外固定一个“绝对形”，就可以用交比重新定义距离和角度，从而得到\textbf{射影度量}。

\begin{definition}
在射影空间中，若额外指定一个固定的二次曲线或二次曲面作为 \textbf{绝对形}，并利用交比来定义点间距离或线间角度，则由此得到的度量结构称为 \textbf{射影度量}。
\end{definition}

这一想法由克莱因系统化提出：欧氏几何、双曲几何、椭圆几何都可以通过“射影几何 + 不同绝对形”的方式得到。

例如，在平面双曲几何的克莱因模型中，绝对形是一条实二次曲线，通常取单位圆。圆内部的点被看作双曲点，圆弦被看作双曲直线，而两点之间的双曲距离可由它们与边界圆的交点构成的交比来定义。这样，交比既可以作为不变量，也可以用来定义“距离”。

对欧氏几何而言，情况更加抽象一些：其绝对形不再是实可见的圆，而是在无穷远直线上一对共轭虚点，也就是所谓的\textbf{虚圆点}。正是借助它们，欧氏角度也可以被重新解释为某种交比的函数。

\begin{remark}
因此，许多几何体系都可以理解成“射影几何 + 某种额外的绝对形”。
\end{remark}

\chapter{二次曲线与配极}

\section{二次曲线的射影观点}
中学解析几何通常把椭圆、抛物线、双曲线分别看成三类不同曲线。但在射影几何中，可以把它们都看成\textbf{二次曲线}，区别只在于它们与无穷远直线的交点情况不同。

\begin{definition}
在射影平面 $\PP^2(\F)$ 中，一个二次齐次方程
\[
Q(x,y,z)=0
\]
的零点集称为一条 \textbf{二次曲线}。
\end{definition}

一般地，二次曲线可写成
\[
Ax^2+Bxy+Cy^2+Dxz+Eyz+Fz^2=0.
\]

如果记
\[
 X=
\begin{pmatrix}
x\\ y\\ z
\end{pmatrix},
\]
那么这条方程还可以写成矩阵形式
\[
 X^{\mathrm T}M X=0,
\]
其中 $M$ 是一个对称矩阵。

这种写法有两个直接的好处：

\begin{enumerate}
    \item 它把几何问题转化为二次型与线性代数问题；
    \item 它便于写出极点极线的定义。
\end{enumerate}

\section{二次曲线的射影分类}
在射影几何里，椭圆、抛物线、双曲线的区别可以概括为：它们与无穷远直线的交点情况不同。

\begin{definition}
设一条实射影平面中的非退化二次曲线与无穷远直线 $\ell_\infty$ 的实交点情况如下：
\begin{enumerate}
    \item 若无实交点，则称为 \textbf{椭圆型}；
    \item 若有一个重实交点（相切），则称为 \textbf{抛物型}；
    \item 若有两个不同实交点，则称为 \textbf{双曲型}。
\end{enumerate}
\end{definition}

分别看三种经典曲线改写成齐次方程后的形式。

\subsection{椭圆型}
椭圆
\[
\frac{x^2}{a^2}+\frac{y^2}{b^2}=1
\]
改写后为
\[
\frac{x^2}{a^2}+\frac{y^2}{b^2}=z^2.
\]
令 $z=0$，得到
\[
\frac{x^2}{a^2}+\frac{y^2}{b^2}=0.
\]
在实数域中，这个方程没有非零解，因此椭圆与无穷远直线没有实交点，故它是椭圆型二次曲线。

\subsection{抛物型}
抛物线
\[
y^2=2px
\]
改写后为
\[
y^2=2pxz.
\]
令 $z=0$，得到
\[
y^2=0.
\]
故在无穷远直线上只有一个重交点 $(1:0:0)$，因此抛物线是抛物型二次曲线。

\subsection{双曲型}
双曲线
\[
\frac{x^2}{a^2}-\frac{y^2}{b^2}=1
\]
改写后为
\[
\frac{x^2}{a^2}-\frac{y^2}{b^2}=z^2.
\]
令 $z=0$，得到
\[
\frac{x^2}{a^2}-\frac{y^2}{b^2}=0,
\]
也即
\[
\frac{x}{a}=\pm \frac{y}{b}.
\]
所以双曲线与无穷远直线有两个实交点，这两个无穷远点正对应双曲线两条渐近线的方向。

\begin{conclusion}
在实射影平面中，区分椭圆、抛物线、双曲线，只要看它们与无穷远直线的实交点个数，分别是 $0,1,2$。
\end{conclusion}

射影变换下三者同属非退化二次曲线；差别体现在它们与无穷远直线的交点情况。

\section{极点与极线}

\subsection{代数定义}
极点极线把二次型中的点和直线配成一对，对偶原理和配极原则都建立在这个点线对应上。

设二次曲线由对称矩阵 $M$ 给出：
\[
 X^{\mathrm T}M X=0.
\]
若取一点
\[
P=[ p],
\]
那么由双线性形式
\[
 p^{\mathrm T}M X=0
\]
所给出的直线，便定义为点 $P$ 的极线。

\begin{definition}
设非退化二次曲线
\[
Q( X)= X^{\mathrm T}M X=0.
\]
对任一点 $P=[ p]$，由方程
\[
 p^{\mathrm T}M X=0
\]
所确定的直线称为点 $P$ 关于该二次曲线的 \textbf{极线}；反过来，给定一条直线，也可唯一对应一个 \textbf{极点}。
\end{definition}

当 $P$ 在曲线上时，该式即切线方程；当 $P$ 不在曲线上时，定义为其极线。

\subsection{仿射形式下的极线方程}
若仿射二次曲线写成
\[
Ax^2+2Bxy+Cy^2+2Dx+2Ey+F=0,
\]
则点 $P(x_0,y_0)$ 的极线方程是
\[
Axx_0+B(xy_0+x_0y)+Cyy_0+D(x+x_0)+E(y+y_0)+F=0.
\]

通常称为“代一半”公式，具体替换如下：

\begin{itemize}
    \item $x^2$ 换成 $xx_0$；
    \item $y^2$ 换成 $yy_0$；
    \item $2xy$ 换成 $xy_0+x_0y$；
    \item $2x$ 换成 $x+x_0$；
    \item $2y$ 换成 $y+y_0$；
    \item 常数项保持不变。
\end{itemize}

\begin{criterion}
对于一般二次曲线
\[
Ax^2+2Bxy+Cy^2+2Dx+2Ey+F=0,
\]
点 $P(x_0,y_0)$ 的极线可用“代一半”快速写成
\[
Axx_0+B(xy_0+x_0y)+Cyy_0+D(x+x_0)+E(y+y_0)+F=0.
\]
\end{criterion}

\subsection{几种常见二次曲线的极线公式}
\subsubsection{圆}
对于圆
\[
x^2+y^2=r^2,
\]
点 $P(x_0,y_0)$ 的极线为
\[
x_0x+y_0y=r^2.
\]

\subsubsection{椭圆}
对于椭圆
\[
\frac{x^2}{a^2}+\frac{y^2}{b^2}=1,
\]
点 $P(x_0,y_0)$ 的极线为
\[
\frac{x_0x}{a^2}+\frac{y_0y}{b^2}=1.
\]

\subsubsection{双曲线}
对于双曲线
\[
\frac{x^2}{a^2}-\frac{y^2}{b^2}=1,
\]
点 $P(x_0,y_0)$ 的极线为
\[
\frac{x_0x}{a^2}-\frac{y_0y}{b^2}=1.
\]

\subsubsection{抛物线}
对于抛物线
\[
y^2=2px,
\]
点 $P(x_0,y_0)$ 的极线为
\[
y_0y=p(x+x_0).
\]

\begin{example}
求椭圆
\[
\frac{x^2}{9}+\frac{y^2}{4}=1
\]
上点 $P(3\cos\theta,2\sin\theta)$ 处的切线方程。
\end{example}

\begin{solution}
由于点 $P$ 在椭圆上，极线退化为该点切线。把 $P$ 的坐标代入椭圆极线公式，得
\[
\frac{(3\cos\theta)x}{9}+\frac{(2\sin\theta)y}{4}=1,
\]
化简为
\[
\frac{x\cos\theta}{3}+\frac{y\sin\theta}{2}=1.
\]
\end{solution}

\subsection{极线的几何意义}
极线的几何意义取决于极点与二次曲线的位置关系：

\begin{enumerate}
    \item 若极点在曲线外，则极线是过该点所作两条切线的切点连线；
    \item 若极点在曲线上，则极线退化为该点处的切线；
    \item 若极点在曲线内，则极线通常与曲线没有实交点，但射影意义下仍完全成立。
\end{enumerate}

所以极线可视为切线方程对平面任意点的推广。

\section{配极原则}
极点极线理论中的一个基本结论是配极原则。它有简洁的代数证明，也有明确的几何意义。

\begin{theorem}
设 $Q=0$ 是一条非退化二次曲线。若点 $P$ 的极线通过点 $Q$，则点 $Q$ 的极线也通过点 $P$。
\end{theorem}

\begin{myproof}
设 $P=[ p]$，$Q=[ q]$。由定义，点 $Q$ 在点 $P$ 的极线上，当且仅当
\[
 p^{\mathrm T}M q=0.
\]
由于矩阵 $M$ 对称，所以
\[
 p^{\mathrm T}M q= q^{\mathrm T}M p.
\]
于是点 $P$ 也在点 $Q$ 的极线上，点线关系的对称性也就成立了。
\end{myproof}

\begin{conclusion}
“$P$ 在 $Q$ 的极线上”与“$Q$ 在 $P$ 的极线上”是完全对称的命题。极点极线对应不是单向的，而是一种自反的点线对偶关系。
\end{conclusion}

\subsection{配极原则的一个重要推论}
如果一点 $P$ 在某条固定直线 $d$ 上运动，那么 $P$ 的极线会恒过某个固定点 $D$。这个点正是直线 $d$ 的极点。

这条结论常用于处理“证明某直线恒过定点”这类问题，即：

\begin{quote}
若一个动点始终落在定直线上，那么它对应的极线恒过该直线的极点。
\end{quote}

同理，若一条动直线始终过某个定点，那么它的极点就会落在某条定直线上。这正说明点线对偶性。

\section{焦点与准线互为极点极线}
这一定理在中学解析几何中通常不会单独提出，但在射影几何中可以自然说明。

\begin{theorem}
对于标准椭圆、双曲线与抛物线，其焦点与对应准线互为极点极线。
\end{theorem}

\begin{myproof}
以椭圆
\[
\frac{x^2}{a^2}+\frac{y^2}{b^2}=1,\qquad c^2=a^2-b^2
\]
为例。右焦点为
\[
F(c,0),
\]
右准线为
\[
x=\frac{a^2}{c}.
\]
把准线写成极线标准形式：
\[
\frac{cx}{a^2}+\frac{0\cdot y}{b^2}=1.
\]
于是对应极点就是 $(c,0)$，即焦点 $F$。故焦点与准线互为极点极线。
\end{myproof}

这条定理说明，焦点准线关系可以看成极点极线在标准坐标下的一个具体写法。

\section{自极三角形}
\begin{definition}
若三角形 $\triangle ABC$ 满足：点 $A$ 的极线是边 $BC$，点 $B$ 的极线是边 $CA$，点 $C$ 的极线是边 $AB$，则称 $\triangle ABC$ 为该二次曲线的 \textbf{自极三角形}。
\end{definition}

在自极三角形里，三组点线对应彼此配合。对每一个顶点而言，它的对边都是自己的极线；反过来，对每一条边而言，它的极点又是对应的顶点。

\begin{property}
自极三角形中的三个顶点两两共轭，三条边也两两对应为共轭线。
\end{property}

自极三角形常用于：

\begin{enumerate}
    \item 化简二次曲线方程；
    \item 写出极点极线；
    \item 组织完全四边形与二次曲线的关系；
    \item 在解析几何中选择更自然的坐标系。
\end{enumerate}

\chapter{对偶原理}

\section{点与直线的对称地位}
射影几何中，点与直线可分别用齐次坐标和齐次线坐标表示，二者具有对偶关系；“点在线上”这一关系本身就是点与直线的双向配合。

\begin{definition}
在射影平面中，将命题中的“点”与“直线”互换，“共线”与“共点”互换，“连接两点”与“求两线交点”互换，若所得命题仍然成立，则称这一现象为 \textbf{对偶原理}。
\end{definition}

对偶原理由射影平面中点、直线的齐次表示自然得到。因为：

\begin{itemize}
    \item 点可写为比例类 $(x:y:z)$；
    \item 直线可写为比例类 $[a:b:c]$；
    \item 点在直线上当且仅当
    \[
    ax+by+cz=0.
    \]
\end{itemize}

这种对称性意味着：点与直线分别属于彼此对偶的空间，而几何命题也常伴随相应的对偶命题。

\section{典型对偶关系}
列出几组常见的对偶对象：

\begin{center}
\begin{tabular}{>{\centering\arraybackslash}m{4cm}>{\centering\arraybackslash}m{4cm}}
\hline
原对象 & 对偶对象 \\
\hline
点 & 直线 \\
共线 & 共点 \\
连接两点成线 & 两线相交得点 \\
点列 & 线束 \\
内接多边形 & 外切多边形 \\
二次点列 & 二级曲线（切线包络） \\
\hline
\end{tabular}
\end{center}

\begin{conclusion}
在射影几何中，一个由点、直线、共点、共线和关联关系叙述的命题，通常都会伴随一个对偶命题。
\end{conclusion}

\section{德萨格定理的自对偶性}
射影几何中的很多定理都具有对偶性，德萨格定理就是一个典型例子。

\begin{theorem}
若两个三角形对应顶点连线共点，则对应边交点共线。
\end{theorem}

把其中的“点”和“直线”互换，“共点”和“共线”互换，就得到：

\begin{theorem}
若两个三角形对应边交点共线，则对应顶点连线共点。
\end{theorem}

这就是德萨格定理的逆命题，也说明德萨格定理具有自对偶性。射影平面内部因此存在稳定的点线对应关系。

\section{对偶图形}
\begin{definition}
在射影平面中，通过点线互换而得到的对应图形称为 \textbf{对偶图形}。
\end{definition}

例如：

\begin{itemize}
    \item 一条直线上的点集构成点列，它的对偶对象是过某点的直线集，即线束；
    \item 一条非退化二次曲线上的点集，其对偶对象是这条曲线的全体切线所形成的包络，也是一条二次对象。
\end{itemize}

特别地，极点极线给出了一种具体的“点 $\leftrightarrow$ 线”对应：它有明确的方程形式，也能转化为题目中的计算条件。

\chapter{几何结构：点列、线束与二次点列}

\section{点列}
\begin{definition}
一条直线上的所有点组成的集合称为 \textbf{点列}。
\end{definition}

点列本身就是一维射影几何的对象，一条射影直线也可以看成一个点列，所以交比、调和点列、对合、射影变换都可以在点列上讨论。

学习射影几何时，可以把点列当作独立的一维对象，不必只把它看成“平面里的特殊情况”。

\section{线束}
\begin{definition}
过同一点的所有直线组成的集合称为 \textbf{线束}。
\end{definition}

线束是点列的对偶对象。过定点的每一条直线都可用方向参数表示，所以线束同样是一维对象；点列上的交比、调和、对合等结论也可以平行地搬到线束上。

例如，若四条共点直线 $a,b,c,d$ 的截线像构成调和点列，那么称这四条线构成调和线束。

\section{点列与线束的透视关系}
一点列可以通过连接到某个中心变成线束；一线束又可以通过与某条截线相交变成点列。因此，点列与线束之间可以互相转化。

点列与线束之间的互相转化，靠的就是“投影—截影”。很多一维结论之所以能在两者之间来回换写，也是因为这种转化不会改变交比等基本量。

\section{二次点列}
\begin{definition}
二次曲线上的点集，在射影几何里称为 \textbf{二次点列}。
\end{definition}

虽然二次点列不在同一直线上，但它在很多方面和一维点列很像。特别是，非退化二次曲线可以通过参数化和射影直线对应起来，因此点列上的不少概念都能搬到二次点列上使用。

例如，可以把二次曲线上四点的交比定义为：任选曲线上一点 $P$，考虑线束
\[
(PA,PB;PC,PD),
\]
并把这个值作为四点 $A,B,C,D$ 的交比。由于它与点 $P$ 的选择无关，这个定义是良定的。

\begin{property}
二次曲线上的四点交比可以借助任何一点对这四点所引射线的交比来定义，因此二次点列同样具有一维射影几何中的交比结构。
\end{property}

因此，二次曲线不仅是一条平面曲线，也常被当作“一维对象”来处理。

\section{射影直线上的复比}
在射影直线 $\PP^1(\F)$ 上，若四个点的齐次坐标分别为
\[
P_1=(u_1:v_1),\quad
P_2=(u_2:v_2),\quad
P_3=(u_3:v_3),\quad
P_4=(u_4:v_4),
\]
则它们的交比可以写成
\[
(P_1,P_2;P_3,P_4)
=
\frac{
\begin{vmatrix}
u_3 & v_3\\
u_1 & v_1
\end{vmatrix}
\begin{vmatrix}
u_4 & v_4\\
u_2 & v_2
\end{vmatrix}
}{
\begin{vmatrix}
u_3 & v_3\\
u_2 & v_2
\end{vmatrix}
\begin{vmatrix}
u_4 & v_4\\
u_1 & v_1
\end{vmatrix}
}.
\]

当取仿射坐标 $x_i=u_i/v_i$ 时，这个式子就退化成熟悉的
\[
(P_1,P_2;P_3,P_4)
=
\frac{(x_3-x_1)(x_4-x_2)}{(x_3-x_2)(x_4-x_1)}.
\]

因此，复比不是某种偶然依赖长度的公式，而是齐次坐标下的一维射影不变量。

\chapter{著名定理}

\section{德萨格定理}
\begin{theorem}
设有两个三角形
\[
\triangle ABC,\qquad \triangle A'B'C'.
\]
若对应顶点连线
\[
AA',\ BB',\ CC'
\]
共点，则对应边交点
\[
AB\cap A'B',\quad BC\cap B'C',\quad CA\cap C'A'
\]
共线。
\end{theorem}

德萨格定理是射影几何中的一个基本定理。它说明“共点透视”与“共线透视”之间存在对应关系。其几何背景也容易说明：若两个三角形位于空间中的不同平面，并且都是由同一个点向这两个平面投影得到，那么它们所在平面的交线恰好就是对应边交点所在的那条直线。

\section{帕普斯定理}
\begin{theorem}
设两条直线 $\ell,m$ 上分别取三点
\[
A,B,C\in \ell,\qquad A',B',C'\in m.
\]
令
\[
X=AB'\cap A'B,\qquad
Y=AC'\cap A'C,\qquad
Z=BC'\cap B'C.
\]
则三点 $X,Y,Z$ 共线。
\end{theorem}

帕普斯定理是早期射影几何中的一个基本点线定理。它表明，仅由两条直线上的六个点和若干交线，就能推出一个共线结论。射影几何中的很多关联结论也可由此得到。

\section{帕斯卡定理}
\begin{theorem}
设六个点
\[
A,B,C,D,E,F
\]
按顺序位于同一条非退化二次曲线上。则三对对边交点
\[
X=AB\cap DE,\qquad
Y=BC\cap EF,\qquad
Z=CD\cap FA
\]
共线。
\end{theorem}

这条共线的直线称为帕斯卡线。帕斯卡定理可以被视为帕普斯定理在二次曲线背景下的推广：当二次曲线退化为两条直线时，帕斯卡定理就会退化成帕普斯定理。

\section{布里昂雄定理}
\begin{theorem}
若一个六边形的六条边都与同一条非退化二次曲线相切，则其三条对顶点连线共点。
\end{theorem}

布里昂雄定理是帕斯卡定理的对偶命题。一个讨论“内接六边形对边交点共线”，另一个讨论“外切六边形对顶点连线共点”。二者共同体现了二次曲线理论中的对偶对称。

\section{射影几何基本定理}
\begin{theorem}
在射影直线 $\PP^1$ 上，给定三个互不相同的点
\[
A,B,C
\]
及另外三个互不相同的点
\[
A',B',C',
\]
则存在唯一的射影变换使得
\[
A\mapsto A',\qquad B\mapsto B',\qquad C\mapsto C'.
\]
\end{theorem}

高维版本的结论是：在 $\PP^n$ 中，一个射影变换由一般位置的 $n+2$ 个点的像决定。因此，射影变换一旦固定足够多的点像，就不再是任意的图形扭曲。

\chapter{射影平面模型}

\section{球面模型}
考虑单位球面
\[
S^2=\{(x,y,z)\in \R^3:x^2+y^2+z^2=1\}.
\]
把球面上任意一对对径点视为同一点，所得空间就可看作实射影平面 $\RP^2$ 的一个模型。

其原因在于：过原点的每一条直线恰好与球面交于一对对径点，而射影点正对应一条过原点的直线。因此，球面模型可以直观看出“射影点是方向类”这一事实。

\section{齐次坐标模型}
齐次坐标模型是常用的代数模型：
\[
\RP^2=\{(x:y:z)\neq (0:0:0)\}/\sim,
\qquad
(x:y:z)\sim(\lambda x:\lambda y:\lambda z),\ \lambda\neq 0.
\]

它的优点在于：

\begin{enumerate}
    \item 点和直线都能写成代数形式；
    \item 射影变换可以直接由矩阵表示；
    \item 二次曲线就是二次齐次方程；
    \item 极点极线可以由二次型与双线性型直接定义。
\end{enumerate}

所以，推导和计算时，齐次坐标模型最方便。

\section{仿射平面加无穷远直线模型}
采用一个常见的模型，把射影平面写成
\[
\RP^2=\A^2\cup \ell_\infty,
\]
其中 $\ell_\infty$ 是无穷远直线。

在这个模型下：

\begin{itemize}
    \item 普通有限点仍然是仿射平面中的点；
    \item 每个方向对应一个无穷远点；
    \item 所有无穷远点组成无穷远直线；
    \item 平行线在无穷远处相交。
\end{itemize}

这个模型对于理解抛物线有一个无穷远点、双曲线有两个无穷远点、椭圆没有实无穷远点很方便。

\section{实射影平面的浸入表示}
实射影平面不能无自交地嵌入三维欧氏空间，但可以浸入。常见的可视化方式是交叉帽（cross-cap）和 Boy 曲面。

要注意的是，虽然克莱因瓶与实射影平面都属于非定向曲面，且都不能在 $\R^3$ 中无自交嵌入，但它们并不是同一种曲面。交叉帽常被用作实射影平面的浸入模型，而克莱因瓶则是另一种独立的非定向闭曲面。

\chapter{内容对应表}

\section{基本依赖关系}
本章概念的使用顺序如下：

\begin{enumerate}
    \item 无穷远点、无穷远直线与齐次坐标给出射影平面的代数表示；
    \item 交比、调和点列、透视对应给出射影变换下保持不变的数量关系；
    \item 射影变换、对偶原理、点列与线束给出点线关系的统一表述；
    \item 二次曲线的射影分类说明椭圆、抛物线、双曲线可视为同一类二次曲线在不同位置下的表现；
    \item 极点极线、配极原则、自极三角形与完全四边形把二次曲线转化为点线对应关系。
\end{enumerate}

\section{常用对应关系}
\begin{center}
\begin{tabular}{@{}ll@{}}
\toprule
几何对象或条件 & 对应的射影表述 \\
\midrule
平行方向 & 无穷远点 \\
平行线族 & 过同一无穷远点的直线束 \\
四点或四线关系 & 交比与调和结构 \\
二次曲线 & 齐次二次型及其退化情形 \\
过定点、在定直线 & 配极关系或对偶关系 \\
切点弦 & 曲线外点的极线 \\
自极三角形 & 三边与三顶点互为极线、极点 \\
\bottomrule
\end{tabular}
\end{center}



